\documentclass[12pt,a4paper]{article}

% ─── Packages ────────────────────────────────────────────────────────────────
\usepackage[T1]{fontenc}
\usepackage{lmodern}
\usepackage{microtype}
\usepackage{geometry}
\geometry{top=1.2in,bottom=1.2in,left=1.3in,right=1.3in}

\usepackage{amsmath,amssymb,amsthm}
\usepackage{mathtools}
\usepackage{tikz-cd}
\usepackage{booktabs}
\usepackage{array}
\usepackage{xcolor}
\usepackage{enumitem}

\usepackage{natbib}

\usepackage{hyperref}
\hypersetup{colorlinks=true,linkcolor=black,citecolor=black,urlcolor=black}

\usepackage{setspace}
\onehalfspacing

% ─── Theorem environments ────────────────────────────────────────────────────
\newtheorem{theorem}{Theorem}[section]
\newtheorem{proposition}[theorem]{Proposition}
\newtheorem{lemma}[theorem]{Lemma}
\newtheorem{corollary}[theorem]{Corollary}
\theoremstyle{definition}
\newtheorem{definition}[theorem]{Definition}
\newtheorem{axiom}{Axiom}
\newtheorem{example}[theorem]{Example}
\theoremstyle{remark}
\newtheorem{remark}[theorem]{Remark}

% ─── Notation ────────────────────────────────────────────────────────────────
\newcommand{\Rnn}{\mathbb{R}_{\ge 0}}
\newcommand{\Ent}{\ensuremath{\mathsf{H}}}
\newcommand{\Opt}{\ensuremath{\mathsf{Opt}}}
\newcommand{\SP}{\mathbf{SP}}
\newcommand{\RSVP}{\mathbf{RSVP}}
\newcommand{\EDSMC}{\mathbf{EDSMC}}
\newcommand{\Ob}{\mathrm{Ob}}
% \newcommand{\Rbi}{\mathbf{R}}
% \newcommand{\Hsym}{\mathsf{H}_{\mathrm{sym}}}
\providecommand{\Hom}{\mathrm{Hom}}
\renewcommand{\Hom}{\mathrm{Hom}}
\providecommand{\id}{\mathrm{id}}
\renewcommand{\id}{\mathrm{id}}
\newcommand{\Pop}{\ensuremath{\mathrm{Pop}}}
\newcommand{\RefOp}{\ensuremath{\mathrm{Ref}}}
\newcommand{\Bind}{\ensuremath{\mathrm{Bind}}}
\newcommand{\Col}{\ensuremath{\mathrm{Col}}}
\newcommand{\Meld}{\mathrm{Meld}}
\newcommand{\tensor}{\otimes}
\newcommand{\unit}{I}
\DeclareMathOperator{\Kc}{K}
\DeclareMathOperator{\colim}{colim}
\newcommand{\PhiF}{\Phi}
\newcommand{\vF}{v}
\newcommand{\SF}{S}
\newcommand{\downset}[1]{\!\downarrow\!{#1}}
\newcommand{\PSh}{\mathbf{PSh}}
\newcommand{\Sh}{\mathbf{Sh}}

% aliases for consistency in the prose
\newcommand{\Refuse}{\RefOp}
\newcommand{\Collapse}{\Col}

% ─── Title ───────────────────────────────────────────────────────────────────

\title{\textbf{Structured Irreversibility}\\
\large Discrete Rewriting and Coherence Realization}

\author{Flyxion}
\date{\today}

\begin{document}

\maketitle

\begin{abstract}
We develop Spherepop as a \emph{free symmetric monoidal entropy-decreasing rewriting
category} $\SP$, prove its initiality in a category $\EDSMC$ of entropy-decreasing
symmetric monoidal categories, and construct a symmetric monoidal functor
$F:\SP\to\RSVP$ into a smooth realization category whose morphisms carry
entropy-slack witnesses ensuring closure.

The first part builds the discrete algebra: option-spaces with admissibility
families, a derived causal preorder, and four irreversible generators---Pop,
Refuse, Bind, and Collapse---each defined set-theoretically. We separate
sequential composition from the symmetric monoidal product and prove the
universal property of Collapse as a coequalizer, the strict symmetric monoidal
structure of $\SP$, and its initiality in $\EDSMC$.

The second part formalizes worldhood via sheaf theory. A presheaf of local
event proposals over an admissible cover satisfies the sheaf condition exactly
when global coherence holds. The meld operator is shown to be policy-induced
sheafification, with decomposition $\Meld_\pi=\mathrm{glue}\circ\Col_\pi$
expressing its universal property.

The third part constructs the realization functor. Objects of $\RSVP$ are
smooth manifolds equipped with coherence potential $\PhiF$, velocity field
$\vF$, and entropy density $\SF$, and morphisms $(\varphi,\eta)$ satisfy
$\varphi^*\SF'=\SF+\eta$. We define $F$ on objects via probability simplices
and on generators via localized smooth perturbations, prove functoriality,
and establish the structural asymmetry: $\SP$ accumulates constraint
(combinatorial, log-bound) while $\RSVP$ redistributes coherence
(differential, field-bound).

We conclude with a worked example, a discrete variational formulation of
commitment, information-theoretic interpretations, normalization results,
and open problems on faithfulness and compact-closed extensions.
\end{abstract}

\tableofcontents
\newpage

% ═════════════════════════════════════════════════════════════════════════════
\section{Introduction}
\label{sec:intro}

\subsection{The Inversion}

Standard computational frameworks take state as primitive and treat time as an index
over state transitions. Spherepop reverses this. The primitive semantic object is not
a state but a \emph{history}: a finite sequence of irreversible events that progressively
restricts an initial space of possibilities. What we call ``state'' is a derived view,
reconstructed by replay. What we call ``meaning'' is incurred through the disciplined
expenditure of possibility, not stored as a representation.

This inversion has both thermodynamic and categorical motivation. From the side of
physics, irreversible computation dissipates degrees of freedom \citep{landauer1961,
bennett1973, GreenAltenkirch2006}: any non-trivial computation in the presence of finite resources must
increase entropy somewhere. From the side of algebra, rewriting systems naturally
generate free categories whose morphisms are reduction sequences. Spherepop combines
these: it treats irreversible rewriting as the fundamental ontological operation and
equips the resulting category with an entropy functional.

\subsection{Structural Quantities}

Before construction, we list the observables and their relationships to prevent
notational confusion later.

\begin{center}
\renewcommand{\arraystretch}{1.4}
\begin{tabular}{lll}
\toprule
\textbf{Quantity} & \textbf{Type} & \textbf{Role} \\
\midrule
Optionality $\Opt(\Omega)$ & $\Ob(\SP)\to\Rnn$ & Degrees of freedom, resource \\
Entropy $\Ent(\Omega)$ & $\Ob(\SP)\to\Rnn$ & Distributional complexity \\
Action $\mathcal{S}[\gamma]$ & histories $\to\Rnn$ & Accumulated optionality cost \\
Commitment $\pi_t$ & step $\to\Rnn$ & Conjugate to optionality \\
Slack $\eta$ & RSVP morphism data & Continuous entropy production \\
Coherence $\PhiF$ & $C^\infty(M)$ & Scalar potential in RSVP \\
\bottomrule
\end{tabular}
\end{center}

Optionality and entropy are distinct observables. Optionality is a pure resource
measure ($\Opt(\Omega)\propto\log|\Omega|$ in the finite uniform case). Entropy
encodes distributional structure and satisfies subadditivity under tensor. Action
is the accumulated optionality cost of a history. Slack is the continuous analogue
of action increments.

\subsection{Two Structures That Must Be Kept Separate}

\paragraph{Composition} is history concatenation:
\[
  (e_2\circ e_1)(\Omega_0) = e_2(e_1(\Omega_0)).
\]
It is strictly associative by construction and encodes temporal order.

\paragraph{Tensor product} $\tensor$ is defined on option-spaces and extended to
morphisms componentwise:
\[
  (\Omega_1,\mathcal{A}_1)\tensor(\Omega_2,\mathcal{A}_2)
  = (\Omega_1\times\Omega_2,\,\mathcal{A}_1\tensor\mathcal{A}_2).
\]
It encodes parallel independence. It is \emph{not} composition.

Entropy behaves differently under each: sequential composition is entropy-monotone
(Axiom~\ref{ax:entropy}); tensor product satisfies subadditivity
(Proposition~\ref{prop:subadditive}).

\subsection{Roadmap}

Section~\ref{sec:ontology} defines option-spaces, admissibility, entropy, and
optionality. Section~\ref{sec:sp} constructs $\SP$ and its generators, with complete
independence proofs. Section~\ref{sec:causality} develops the causal preorder and
proves the full universal property of Collapse. Section~\ref{sec:monoidal} establishes
the strict symmetric monoidal structure and proves initiality of $\SP$ in $\EDSMC$.
Section~\ref{sec:sheaf} formalizes worldhood and proves the meld theorem.
Section~\ref{sec:rsvp} constructs $\RSVP$ with witnessed morphisms.
Section~\ref{sec:functor} defines $F$ and proves its four component lemmas and
overall functoriality. Section~\ref{sec:worked} works through a complete example.
Later sections address mechanics, information theory, normalization, further
directions, and failure modes.

% ═════════════════════════════════════════════════════════════════════════════
\section{Option-Spaces, Admissibility, and Entropy}
\label{sec:ontology}

\subsection{Option-Spaces}

\begin{definition}[Option-space]
\label{def:optspace}
An \emph{option-space} is a pair $(\Omega,\mathcal{A})$ where:
\begin{enumerate}[label=(\roman*)]
\item $\Omega$ is a nonempty set of \emph{admissible futures};
\item $\mathcal{A}\subseteq 2^\Omega$ is the \emph{admissibility family}, satisfying
  $\Omega\in\mathcal{A}$ and closure under finite intersection:
  $U,V\in\mathcal{A}\Rightarrow U\cap V\in\mathcal{A}$.
\end{enumerate}
\end{definition}

The family $\mathcal{A}$ encodes structural compatibility constraints imposed by
binding events. A subset $U\subseteq\Omega$ lies in $\mathcal{A}$ precisely when it
respects all currently active dependency constraints.

\begin{example}
Let $\Omega=\{a,b,c\}$ and $\mathcal{A}=\{\Omega,\{a,b\},\{a,c\},\{a\}\}$. Then
$a$ is causally necessary for both $b$ and $c$: no admissible future can contain $b$
or $c$ without also containing $a$.
\end{example}

\subsection{Entropy and Optionality}

\begin{definition}[Entropy functional]
\label{def:entropy}
A \emph{Spherepop entropy functional} is a map $\Ent:\Ob(\SP)\to\Rnn$ satisfying:
\begin{enumerate}[label=(\roman*)]
\item $\Ent(\Omega)=0$ iff $|\Omega|=1$;
\item \emph{Monotonicity:} $\Omega'\subseteq\Omega$ implies $\Ent(\Omega')\le\Ent(\Omega)$;
\item \emph{Subadditivity:} $\Ent(\Omega_1\tensor\Omega_2)\le\Ent(\Omega_1)+\Ent(\Omega_2)$,
  with equality when $\Omega_1,\Omega_2$ are independent.
\end{enumerate}
\end{definition}

The canonical instance: for a probability measure $p$ on $\Omega$,
$\Ent(\Omega)=-\sum_{x\in\Omega}p(x)\log p(x)$.

\begin{definition}[Optionality]
\label{def:opt}
The \emph{optionality functional} is a monotone map $\Opt:\Ob(\SP)\to\Rnn$
satisfying $\Omega'\subseteq\Omega\Rightarrow\Opt(\Omega')\le\Opt(\Omega)$ and
$\Opt(\{\omega\})=0$.
\end{definition}

Optionality measures degrees of freedom available to a world. In the finite uniform
case one may take $\Opt(\Omega)=\log|\Omega|=\Ent(\Omega)$; in general they differ.

\paragraph{Summary.} The ontological layer is now fixed: option-spaces are the objects;
entropy measures distributional complexity; optionality measures resource availability.
These are the inputs to the algebraic construction of $\SP$.

\section{Entropy as a Monoidal Weight} \label{sec:entropy-weight}

The entropy functional $\Ent : \Ob(\SP) \to \Rnn$ has so far been treated as a monotone observable. We now strengthen its role: entropy induces an enrichment of $\SP$ over the poset $(\Rnn, \le)$ and yields a strict monoidal weight structure.
\subsection{Entropy-Decreasing Enrichment}

For objects $\Omega,\Omega'\in\Ob(\SP)$, define a preorder on $\Hom_{\SP}(\Omega,\Omega')$ by
\[
e \preceq e' \quad \Longleftrightarrow \quad \Ent(e(\Omega)) \ge \Ent(e'(\Omega)).
\]
Compatibility with composition is expressed by the monotonicity rule
\begin{align*}
e_1 \preceq e_1' \ \text{and}\ e_2 \preceq e_2'
\quad \Longrightarrow \quad
e_2\circ e_1 \preceq e_2'\circ e_1'.
\end{align*}

\begin{proposition}
$\SP$ is enriched over the monoidal poset $(\Rnn,\ge,+,0)$, with enrichment given by entropy decrease.
\end{proposition}

\begin{proof}
Define the weight of $e:\Omega\to\Omega'$ to be
\[
w(e) := \Ent(\Omega)-\Ent(\Omega').
\]
Then for composable morphisms $\Omega_0\xrightarrow{e_1}\Omega_1\xrightarrow{e_2}\Omega_2$ we have
\[
w(e_2\circ e_1)=\Ent(\Omega_0)-\Ent(\Omega_2)=\bigl(\Ent(\Omega_0)-\Ent(\Omega_1)\bigr)+\bigl(\Ent(\Omega_1)-\Ent(\Omega_2)\bigr)=w(e_1)+w(e_2),
\]
which is strict additivity. The remaining enrichment axioms follow from the definition of $\preceq$ and Axiom~\ref{ax:entropy}.
\end{proof}

This makes $\SP$ a \emph{strictly additive entropy-weighted category}. In particular, histories admit a canonical factorization by total entropy drop.

\subsection{Additive Decomposition of Histories}

Any history $h$ admits a decomposition of the schematic form
\[
h \;=\; \Collapse^{k_c}\circ \Bind^{k_b}\circ \Pop^{k_p},
\]
and the total weight decomposes additively as
\[
w(h)
=\sum_{i=1}^{k_p} w(\Pop_i) + \sum_{j=1}^{k_c} w(\Collapse_j),
\]
with the understanding that $\Bind$-steps contribute zero Shannon-entropy weight in the current convention.


\section{Bi-Graded Resource Functional}
\label{sec:bigraded}

The entropy functional $\Ent$ of Section~\ref{sec:ontology} is Shannon-type and
measures distributional complexity. As shown in Appendix~\ref{app:symm-entropy},
there is a second natural observable, $\Hsym(\Omega,\mathcal{A})=\log|\Omega|
-\log|\mathrm{Aut}(\Omega,\mathcal{A})|$, which captures \emph{symmetry entropy}:
the log-ratio of raw size to structural symmetry. The two functionals behave
differently across generators—$\Ent$ is strictly decreased by $\Pop$ and $\RefOp$
and preserved by $\Bind$, while $\Hsym$ is strictly increased by $\Bind$. This
section consolidates them into a single bi-graded resource functional and shows
that the result is both a Lawvere enrichment and a natural basis for the
factorization system of Section~\ref{sec:factorization}.

\subsection{The Bi-Graded Ordered Monoid}

\begin{definition}[Bi-graded resource monoid]
\label{def:bigraded-monoid}
Let $\mathcal{R}:=([0,\infty]\times[0,\infty],\ge_{\times},+)$ where:
\begin{itemize}
  \item addition is componentwise: $(a,b)+(a',b'):=(a+a',b+b')$;
  \item the partial order is componentwise: $(a,b)\ge_\times(a',b')$ iff
    $a\ge a'$ and $b\ge b'$.
\end{itemize}
$(\mathcal{R},\ge_\times,+,(0,0))$ is a commutative ordered monoid, hence a
symmetric monoidal category when regarded as thin.
\end{definition}

\begin{definition}[Bi-graded resource functional]
\label{def:bigraded-resource}
Define $\mathbf{R}:\Ob(\SP)\to\mathcal{R}$ by
\[
  \mathbf{R}(\Omega,\mathcal{A})
  :=\bigl(\Ent(\Omega,\mathcal{A}),\;\Hsym(\Omega,\mathcal{A})\bigr),
\]
where $\Ent$ is the Shannon entropy functional
(Definition~\ref{def:entropy}) and $\Hsym(\Omega,\mathcal{A}):=
\log|\Omega|-\log|\mathrm{Aut}(\Omega,\mathcal{A})|$ is the symmetry entropy
(Appendix~\ref{app:symm-entropy}).
\end{definition}

\subsection{Generator Behavior Under the Bi-Graded Functional}

The following table records the direction of change in each component under each
generator. Let $\downarrow_s$ denote strict decrease, $\uparrow_s$ strict increase,
and $=$ no change.

\begin{center}
\renewcommand{\arraystretch}{1.4}
\begin{tabular}{lll}
\toprule
\textbf{Generator} & $\Ent$ component & $\Hsym$ component \\
\midrule
$\Pop_U$ & $\downarrow_s$ & depends on automorphisms of $\Omega\setminus U$ \\
$\RefOp_U$ & $\downarrow_s$ & depends on automorphisms of $\Omega\setminus U$ \\
$\Bind_{ij}$ & $=$ & $\downarrow_s$ (breaks the $i\leftrightarrow j$ automorphism) \\
$\Col_\sim$ & $\downarrow_s$ or $=$ & $\downarrow_s$ or $=$ \\
\bottomrule
\end{tabular}
\end{center}

\begin{proposition}[$\mathbf{R}$ is resource-monotone]
\label{prop:bigraded-mono}
For every morphism $f:(\Omega,\mathcal{A})\to(\Omega',\mathcal{A}')$ in $\SP$,
\[
  \mathbf{R}(\Omega',\mathcal{A}')\le_\times\mathbf{R}(\Omega,\mathcal{A}),
\]
with strict decrease in at least one component for every non-identity morphism.
\end{proposition}

\begin{proof}
For each generator:
\begin{itemize}
  \item $\Pop_U$, $\RefOp_U$: $\Ent$ strictly decreases (Axiom~\ref{ax:entropy}).
    The $\Hsym$ component may or may not change, but cannot increase past the
    original value since $|\Omega\setminus U|<|\Omega|$ bounds $\log|\Omega\setminus U|$.
  \item $\Bind_{ij}$: $\Ent$ is preserved (Axiom~\ref{ax:entropy}). The automorphism
    that swaps $i\leftrightarrow j$ is removed from $\mathrm{Aut}(\Omega,\mathcal{A}')$
    (Appendix~\ref{app:symm-entropy}), so $|\mathrm{Aut}(\Omega,\mathcal{A}')|<
    |\mathrm{Aut}(\Omega,\mathcal{A})|$ and $\Hsym$ strictly decreases.
  \item $\Col_\sim$: $\Ent$ decreases or stays equal. The coarse-graining reduces
    $|\Omega/{\sim}|\le|\Omega|$ and may reduce automorphisms; $\Hsym$ does not
    increase.
\end{itemize}
In each case, $\mathbf{R}$ is non-increasing componentwise, with a strict decrease
in at least one component.
\end{proof}

\subsection{Bi-Graded Lawvere Enrichment}

The bi-graded functional extends the Lawvere enrichment of Section~\ref{sec:lawvere}
to two dimensions.

\begin{definition}[Bi-graded distance]
\label{def:bigraded-dist}
Define $d_{\mathbf{R}}:\Ob(\SP)\times\Ob(\SP)\to\mathcal{R}$ by
\[
  d_{\mathbf{R}}(X,Y)
  :=\inf_{\ge_\times}\bigl\{\mathbf{R}(Y)-\mathbf{R}(X):\exists f:Y\to X\bigr\},
\]
where the infimum is taken componentwise and $\mathbf{R}(Y)-\mathbf{R}(X)$ is
defined componentwise when all entries are finite.
\end{definition}

\begin{proposition}
$(\Ob(\SP),d_{\mathbf{R}})$ is a $\mathcal{R}$-enriched category (a generalized
Lawvere metric space valued in $\mathcal{R}$).
\end{proposition}

\begin{proof}
$d_{\mathbf{R}}(X,X)=(0,0)$ (use $\id_X$). Triangle inequality: the argument of
Proposition~\ref{prop:lawvere-metric} applies componentwise.
\end{proof}

\subsection{Refined Factorization System}

The bi-graded functional sharpens the $\mathcal{E}$--$\mathcal{M}$ factorization of
Section~\ref{sec:factorization}.

\begin{proposition}[Bi-graded factorization]
\label{prop:bigraded-factor}
Under $\mathbf{R}$:
\begin{itemize}
  \item $\mathcal{E}$ (the $\Bind$-generated morphisms) are exactly the morphisms
    that move only the $\Hsym$ component: $\Delta\Ent=0$, $\Delta\Hsym<0$.
  \item $\mathcal{M}$ (the $\Pop$/$\RefOp$/$\Col$-generated morphisms) are exactly
    the morphisms that move the $\Ent$ component: $\Delta\Ent<0$ (or $=0$ for
    entropy-neutral $\Col$), with $\Delta\Hsym$ unconstrained.
\end{itemize}
Hence the $\mathcal{E}$--$\mathcal{M}$ factorization of every morphism in $\SP$ is
a factorization into a \emph{symmetry-reducing} step followed by an
\emph{option-eliminating} step.
\end{proposition}

\begin{proof}
Immediate from Proposition~\ref{prop:bigraded-mono} and the generator analysis above:
bindings move only $\Hsym$; pops, refuses, and collapses move $\Ent$. The factorization
theorem of Section~\ref{sec:factorization} (commuting $\Bind$ steps past $\Pop/\Col$
when supports are disjoint) expresses every history in this two-phase form.
\end{proof}

\begin{remark}
The bi-graded resource $\mathbf{R}$ provides a principled resolution to a subtlety
in the entropy-alone treatment: an entropy-neutral $\Col$ (one that identifies
entropy-equal classes) contributes zero to $\Ent$ but still decreases $\Hsym$,
so it appears as a non-trivial morphism under $\mathbf{R}$. The bi-grading therefore
distinguishes between identity morphisms and genuinely non-trivial entropy-neutral
collapses, a distinction invisible to the single-component functional.
\end{remark}

\section{Lawvere-Metric Enrichment of \texorpdfstring{$\SP$}{SP}}
\label{sec:lawvere}

Section~\ref{sec:entropy-weight} showed that entropy differences compose
additively along morphism chains. We now consolidate this into a formal
enrichment over the Lawvere quantale, giving the entropy grading the status
of a genuine categorical metric structure.

\subsection{The Lawvere Quantale}

Recall that the \emph{Lawvere quantale} is the monoidal poset
$\mathbf{L}:=([0,\infty],\ge,+,0)$, regarded as a thin monoidal category:
there is a unique morphism $a\to b$ in $\mathbf{L}$ if and only if $a\ge b$,
composition is addition, and the unit object is $0$.
An \emph{$\mathbf{L}$-enriched category} (equivalently, a generalized metric
space in the sense of Lawvere \citep{lawvere1970}) is a set of objects $X$
together with a function $d:X\times X\to[0,\infty]$ satisfying
\[
  d(X,X)=0 \qquad\text{and}\qquad d(X,Z)\le d(X,Y)+d(Y,Z).
\]
Symmetry is \emph{not} required; this directedness is essential for capturing
the asymmetry of irreversible entropy flow.

\subsection{The Entropy Distance on \texorpdfstring{$\Ob(\SP)$}{Ob(SP)}}

\begin{definition}[Entropy distance]
\label{def:entropy-dist}
Define $d_{\Ent}:\Ob(\SP)\times\Ob(\SP)\to[0,\infty]$ by
\[
  d_{\Ent}(X,Y) :=
  \inf\bigl\{\Ent(Y)-\Ent(X)\;:\;\exists f:Y\to X\text{ in }\SP\bigr\},
\]
with the convention $\inf\emptyset=\infty$.
\end{definition}

The direction of the morphism $f:Y\to X$ is intentional: by the entropy
preorder of Section~\ref{sec:structured}, $X\preceq Y$ iff such an $f$
exists, and Axiom~\ref{ax:entropy} then forces $\Ent(X)\le\Ent(Y)$, making
every term in the infimum non-negative.

\begin{proposition}[$d_{\Ent}$ is a Lawvere metric]
\label{prop:lawvere-metric}
$(\Ob(\SP),d_{\Ent})$ is a $\mathbf{L}$-enriched category. Explicitly:
\begin{enumerate}[label=(\roman*)]
  \item $d_{\Ent}(X,X)=0$ for all $X$.
  \item $d_{\Ent}(X,Z)\le d_{\Ent}(X,Y)+d_{\Ent}(Y,Z)$ for all $X,Y,Z$.
\end{enumerate}
\end{proposition}

\begin{proof}
\textbf{(i).} The identity morphism $\id_X:X\to X$ witnesses $\Ent(X)-\Ent(X)=0$,
so the infimum is at most $0$; it is at least $0$ by Axiom~\ref{ax:entropy}.

\textbf{(ii).} Fix $\varepsilon>0$. Choose $f:Y\to X$ with
$\Ent(Y)-\Ent(X)<d_{\Ent}(X,Y)+\varepsilon$, and $g:Z\to Y$ with
$\Ent(Z)-\Ent(Y)<d_{\Ent}(Y,Z)+\varepsilon$. The composite $f\circ g:Z\to X$
is a morphism in $\SP$, so
\[
  d_{\Ent}(X,Z)\le\Ent(Z)-\Ent(X)
  =\bigl(\Ent(Z)-\Ent(Y)\bigr)+\bigl(\Ent(Y)-\Ent(X)\bigr)
  <d_{\Ent}(X,Y)+d_{\Ent}(Y,Z)+2\varepsilon.
\]
Since $\varepsilon>0$ was arbitrary, the triangle inequality holds.
\end{proof}

\subsection{Collapse to a Clean Formula}

When $\Ent$ is strictly functorial (each generator contributes a determined entropy
drop and these drops telescope), the infimum in Definition~\ref{def:entropy-dist}
is attained and simplifies completely.

\begin{proposition}[Telescoping collapse]
\label{prop:telescoping}
For any composable history $\gamma:Y\to X$, the path length
$\ell_{\Ent}(\gamma):=\sum_i\delta_{\Ent}(e_i)$ equals $\Ent(Y)-\Ent(X)$.
Consequently, if $X\preceq Y$ (i.e.\ $\exists f:Y\to X$), then
\[
  d_{\Ent}(X,Y) = \Ent(Y)-\Ent(X).
\]
If no such morphism exists, $d_{\Ent}(X,Y)=\infty$.
\end{proposition}

\begin{proof}
Telescoping of the entropy functional along the history, as in the proof of the
weight-additivity Proposition in Section~\ref{sec:entropy-weight}. The formula
for $d_{\Ent}$ then follows because every path $Y\to X$ has the same length
$\Ent(Y)-\Ent(X)$, so the infimum is that common value.
\end{proof}

\begin{remark}
Proposition~\ref{prop:telescoping} upgrades the Lawvere enrichment to the
strongest possible form: $d_{\Ent}$ is the \emph{canonical} Lawvere metric
on the entropy-preordered set of objects, with distance equal to the entropy
gap when reachable and $\infty$ otherwise. The asymmetry
$d_{\Ent}(X,Y)\ne d_{\Ent}(Y,X)$ in general (since $X\preceq Y$ does not
imply $Y\preceq X$) reflects the irreversibility of $\SP$ at the
metrized level.
\end{remark}

\subsection{Thin-Category Formulation}

For the purposes of enriched category theory, it is cleaner to record the
enrichment on hom-objects rather than on object-pairs. Define the
\emph{thin reachability category} $\Pi(\SP)$: same objects as $\SP$;
$\Hom_{\Pi(\SP)}(Y,X)$ is a singleton $\{\ast\}$ if $\SP(Y,X)\ne\emptyset$,
and empty otherwise. Equip $\Pi(\SP)$ with the $\mathbf{L}$-hom-object
\[
  \underline{\mathrm{Hom}}(Y,X):=d_{\Ent}(X,Y)\in[0,\infty].
\]
Composition in $\Pi(\SP)$ sends $(d_1,d_2)\mapsto d_1+d_2$, which is
precisely the $\mathbf{L}$-composition. This makes $\Pi(\SP)$ a
$\mathbf{L}$-enriched category in the full sense of \citet{awodey2010},
and $d_{\Ent}$ is its hom-functor.

\section{A Canonical Factorization System} \label{sec:factorization}

The four generators admit a canonical orthogonal factorization system separating \emph{structural introduction} from \emph{option elimination}.

\subsection{E--M Factorization}

Define two classes of morphisms:
\[
\mathcal{E} := \{\id,\ \Bind\text{-generated morphisms}\},
\qquad
\mathcal{M} := \{\Pop,\ \RefOp,\ \Col\text{-generated morphisms}\}.
\]

\begin{theorem}
Every morphism $h$ in $\SP$ factors uniquely (up to policy equivalence) as
\[
h = m \circ e
\]
with $e\in\mathcal{E}$ and $m\in\mathcal{M}$.
\end{theorem}

\begin{proof}[Sketch]
Bindings refine admissibility families while leaving the underlying set $\Omega$ unchanged, hence they preserve the Shannon-entropy functional under the convention of Axiom~\ref{ax:entropy}. Pops and collapses are the only generators that change the underlying option-set and therefore the only generators that can contribute to strict entropy descent. The factorization is obtained by commuting independent $\Bind$ steps leftward past $\Pop/\RefOp/\Col$ steps whenever their supports are disjoint, using the commutation relations already imposed in Definition~\ref{def:free-gen}. Uniqueness holds up to the same policy congruence used to identify commuting independent bindings.
\end{proof}

% ═════════════════════════════════════════════════════════════════════════════
\section{The Category \texorpdfstring{$\SP$}{SP}}
\label{sec:sp}

\subsection{Objects, Morphisms, and Axioms}

Objects of $\SP$ are option-spaces. A morphism $e:(\Omega,\mathcal{A})\to
(\Omega',\mathcal{A}')$ is an \emph{irreversible event}: a function $e:\Omega\to\Omega'$
that is admissibility-preserving ($A'\in\mathcal{A}'\Rightarrow e^{-1}(A')\in\mathcal{A}$)
and entropy-monotone.

\begin{axiom}[Sequential irreversibility]
\label{ax:irrev}
No non-identity morphism in $\SP$ admits an inverse.
\end{axiom}

\begin{axiom}[Entropy monotonicity]
\label{ax:entropy}
For every morphism $e:\Omega\to\Omega'$: $\Ent(\Omega')\le\Ent(\Omega)$.
Strict decrease holds for $\Pop$ and $\RefOp$. $\Bind$ preserves entropy while
reducing symmetry. $\Col$ may produce equality or strict decrease depending on
whether the quotient collapses entropy-distinct classes.
\end{axiom}

\subsection{Generating Morphisms}

\begin{definition}[Pop]
\label{def:pop}
A $\Pop$-event eliminates a nonempty $U\subsetneq\Omega$:
\[
  \Pop_U : (\Omega,\mathcal{A}) \longrightarrow
  (\Omega\setminus U,\;\{A\cap(\Omega\setminus U):A\in\mathcal{A}\}).
\]
Entropy decreases strictly: $\Ent(\Omega\setminus U)<\Ent(\Omega)$.
\end{definition}

\begin{definition}[Refuse]
\label{def:ref}
A $\RefOp$-event has the same underlying set-theoretic action as $\Pop_U$ but carries
an auxiliary tag $\tau\in\mathcal{B}(\Omega)$ recorded by an accounting functor
$\mathcal{B}:\SP\to\mathbf{Set}$, distinguishing deliberate non-selection from
positive elimination. The underlying morphism of option-spaces coincides with $\Pop_U$.
\end{definition}

\begin{definition}[Bind]
\label{def:bind}
A $\Bind$-event tightens the admissibility family without eliminating futures:
\[
  \Bind_{ij} : (\Omega,\mathcal{A}) \longrightarrow
  (\Omega,\,\mathcal{A}')
\]
where $\mathcal{A}'=\{A\in\mathcal{A}: i\in A\Rightarrow j\in A\}\subsetneq\mathcal{A}$.
Entropy is globally preserved: $\Ent(\Omega,\mathcal{A}')=\Ent(\Omega,\mathcal{A})$.
\end{definition}

\begin{definition}[Collapse]
\label{def:col}
A $\Col$-event applies a causally admissible equivalence relation $\sim$ (see
Definition~\ref{def:causal-admissible}) to form the quotient:
\[
  \Col_\sim : (\Omega,\mathcal{A}) \longrightarrow
  (\Omega/{\sim},\,\mathcal{A}/{\sim}).
\]
\end{definition}

\subsection{Free Generation and Independence of Generators}

\begin{definition}[Free generation]
\label{def:free-gen}
Let $\mathcal{G}$ be the directed multigraph with vertices $\Ob(\SP)$ and edges
labeled by instances of the four generators. $\SP$ is the quotient of the free
category $\mathcal{F}(\mathcal{G})$ by the congruence generated by:
\begin{enumerate}[label=(\alph*)]
\item entropy monotonicity (Axiom~\ref{ax:entropy});
\item idempotence of Pop on already-excluded elements: $\Pop_U\circ\Pop_V
  =\Pop_{U\cup V}$ when $U\cap(\Omega\setminus V)=\emptyset$;
\item commutativity of independent Bind events: $\Bind_{ij}\circ\Bind_{kl}
  =\Bind_{kl}\circ\Bind_{ij}$ when $\{i,j\}\cap\{k,l\}=\emptyset$;
\item the symmetric monoidal laws (Section~\ref{sec:monoidal}).
\end{enumerate}
No further relations are imposed.
\end{definition}

\begin{proposition}[Independence of generators]
\label{prop:independence}
The four generators are independent in the following precise sense:
\begin{enumerate}[label=(\roman*)]
\item $\Pop$ cannot be derived from $\RefOp$: they differ in the value of the
  accounting functor $\mathcal{B}$.
\item $\RefOp$ cannot be derived from $\Pop$: no composition of $\Pop$ events
  yields a morphism with a non-trivial tag in $\mathcal{B}$.
\item $\Bind$ is not reducible to $\Pop$ followed by $\Col$: $\Bind$ preserves
  $|\Omega|$ while any $\Pop$ reduces it; a subsequent $\Col$ could re-identify
  elements but cannot restore the eliminated element without violating
  Axiom~\ref{ax:irrev}.
\item $\Col$ is not a special case of $\Pop$: $\Pop$ produces a strict subobject;
  $\Col$ produces a quotient. These are categorically distinct constructions.
\end{enumerate}
\end{proposition}

\begin{proof}
For (i) and (ii): the accounting functor $\mathcal{B}$ assigns distinct values to
$\Pop$ and $\RefOp$ events by definition. Since $\mathcal{B}$ is a functor, it must
agree on composed morphisms with the composed images; but $\Pop$ has $\mathcal{B}$
value $0$ and $\RefOp$ has value $\tau\ne 0$, so neither can be derived from a
composition involving only the other.

For (iii): let $e$ be any morphism derivable from $\Pop$ and $\Col$ events starting
at $(\Omega,\mathcal{A})$. The $\Pop$ step yields $(\Omega\setminus U,\cdot)$ with
$|\Omega\setminus U|<|\Omega|$. A subsequent $\Col$ reduces $|\Omega\setminus U|$
further or preserves it. The resulting object has $|\Omega''|\le|\Omega\setminus U|
<|\Omega|$. But $\Bind_{ij}$ produces $(\Omega,\mathcal{A}')$ with $|\Omega|$
unchanged. Hence $\Bind$ morphisms lie outside the image of the subcategory generated
by $\Pop$ and $\Col$.

For (iv): $\Pop_U$ is the inclusion $\Omega\setminus U\hookrightarrow\Omega$ reversed
(a restriction), hence an injective-on-objects construction. $\Col_\sim$ is a
surjective-on-objects quotient. In a category of structured sets, injective and
surjective constructions on objects are distinct.
\end{proof}

\paragraph{Summary.} The category $\SP$ is freely generated by the four independent
generators modulo the relations of Definition~\ref{def:free-gen}. Composition encodes
temporal sequence. No generator is reducible to any combination of the others.
\subsection{Entropy Slack Witnesses}

To ensure strict functoriality of $F$, we refine RSVP morphisms.

\begin{definition}
An RSVP morphism is a pair $(\varphi,\eta)$ where
\[
\varphi:M\to M', \qquad \eta\in C^\infty(M,\Rnn), \qquad \varphi^*\SF'=\SF+\eta.
\]
The function $\eta$ is the \emph{entropy slack witness}. Composition is defined by
\[
(\varphi_2,\eta_2)\circ(\varphi_1,\eta_1)
=
(\varphi_2\circ\varphi_1,\; \eta_1+\varphi_1^*\eta_2).
\]
\end{definition}

\begin{proposition}
With slack witnesses, $\RSVP$ is strictly closed under composition.
\end{proposition}

\begin{proof}
Pullback preserves addition, so
\[
(\varphi_2\circ\varphi_1)^*\SF''
=\varphi_1^*(\varphi_2^*\SF'')
=\varphi_1^*(\SF'+\eta_2)
=\varphi_1^*\SF'+\varphi_1^*\eta_2
=(\SF+\eta_1)+\varphi_1^*\eta_2
=\SF+(\eta_1+\varphi_1^*\eta_2),
\]
which is exactly the defining condition for the composite slack.
\end{proof}

% ═════════════════════════════════════════════════════════════════════════════
\section{Causality and the Universal Property of Collapse}
\label{sec:causality}

\subsection{Causal Preorder}

\begin{definition}[Causal preorder]
\label{def:causal-preorder}
For $(\Omega,\mathcal{A})$ and $x,y\in\Omega$, define $x\preceq y$ iff every
$A\in\mathcal{A}$ containing $y$ also contains $x$:
\[
  x\preceq y \;\iff\; \forall A\in\mathcal{A}:\; y\in A\Rightarrow x\in A.
\]
Write $\downset{y}=\{x\in\Omega:x\preceq y\}$ for the causal past of $y$.
\end{definition}

\begin{lemma}[Reflexivity and transitivity]
\label{lem:preorder}
$\preceq$ is a preorder.
\end{lemma}

\begin{proof}
\textit{Reflexivity.} For any $x$: every $A$ containing $x$ contains $x$. Hence
$x\preceq x$.

\textit{Transitivity.} Suppose $x\preceq y$ and $y\preceq z$. Let $A\in\mathcal{A}$
with $z\in A$. Since $y\preceq z$, we have $y\in A$. Since $x\preceq y$, we have
$x\in A$. Hence $x\preceq z$.
\end{proof}

\begin{lemma}[Antisymmetry under acyclicity]
\label{lem:antisymmetry}
If the admissibility family $\mathcal{A}$ contains no nontrivial cycles under
$\preceq$ (i.e. $x\preceq y$ and $y\preceq x$ implies $x=y$), then $\preceq$ is a
partial order.
\end{lemma}

\begin{proof}
The hypothesis directly gives antisymmetry. Combined with Lemma~\ref{lem:preorder},
$\preceq$ is a partial order.
\end{proof}

\begin{remark}
$\Bind_{ij}$ adds the constraint $i\preceq j$ to $\mathcal{A}$. Acyclicity of the
dependency graph (no sequence of $\Bind$ events creates a cycle $i\preceq j\preceq i$
with $i\ne j$) is a well-formedness condition on histories. We assume it throughout.
\end{remark}

\subsection{Causally Admissible Equivalence Relations}

\begin{definition}[Causally admissible equivalence]
\label{def:causal-admissible}
An equivalence relation $\sim$ on $\Omega$ is \emph{causally admissible} if
\[
  x\sim y \;\implies\; \downset{x}=\downset{y}.
\]
\end{definition}

This condition ensures that identified futures have identical causal pasts, so
collapse does not alter downstream commitment structure.

\subsection{Full Universal Property of Collapse}

Fix $(\Omega,\mathcal{A})$ and a causally admissible equivalence $\sim_q$. Let
$\pi_q:\Omega\twoheadrightarrow\Omega/{\sim_q}$ be the quotient map and
$\mathcal{A}_q=\{A\subseteq\Omega/{\sim_q}:\pi_q^{-1}(A)\in\mathcal{A}\}$ the
induced admissibility family. Let $\SP_c$ denote the full subcategory of $\SP$
of causal-preorder-preserving morphisms.

\begin{theorem}[Universal property of causal collapse]
\label{thm:collapse-universal}
In $\SP_c$, the morphism $\Col_q:(\Omega,\mathcal{A})\to(\Omega/{\sim_q},\mathcal{A}_q)$
is a coequalizer of the kernel pair of $\pi_q$. Explicitly: define
$R_q=\{(x,y)\in\Omega^2:x\sim_q y\}$ with induced admissibility, and projections
$r_1,r_2:(R_q,\mathcal{A}_{R_q})\rightrightarrows(\Omega,\mathcal{A})$,
$r_1(x,y)=x$, $r_2(x,y)=y$. Then:

\medskip
\noindent\textbf{(i) Coequalizing.}
$\Col_q\circ r_1=\Col_q\circ r_2$.

\medskip
\noindent\textbf{(ii) Universal property.}
For any $(\Omega'',\mathcal{A}'')\in\SP_c$ and any morphism
$f:(\Omega,\mathcal{A})\to(\Omega'',\mathcal{A}'')$ in $\SP_c$ with
$f\circ r_1=f\circ r_2$, there exists a unique morphism
$\bar{f}:(\Omega/{\sim_q},\mathcal{A}_q)\to(\Omega'',\mathcal{A}'')$ in $\SP_c$
with $\bar{f}\circ\Col_q=f$:
\[
\begin{tikzcd}[column sep=3.5em, row sep=2.5em]
  (R_q,\mathcal{A}_{R_q})
    \arrow[r, shift left=0.3em, "r_1"]
    \arrow[r, shift right=0.3em, "r_2"']
  &
  (\Omega,\mathcal{A})
    \arrow[r, "\Col_q"]
    \arrow[dr, "f"']
  &
  (\Omega/{\sim_q},\mathcal{A}_q)
    \arrow[d, dashed, "\exists!\,\bar{f}"]
  \\
  & &
  (\Omega'',\mathcal{A}'')
\end{tikzcd}
\]

\medskip
\noindent\textbf{(iii) Naturality.}
The assignment $f\mapsto\bar{f}$ is natural in $(\Omega'',\mathcal{A}'')$.
\end{theorem}

\begin{proof}
\textbf{Part (i).} For $(x,y)\in R_q$, we have $x\sim_q y$, hence
$\pi_q(x)=\pi_q(y)$. Thus $\Col_q(r_1(x,y))=\pi_q(x)=\pi_q(y)=\Col_q(r_2(x,y))$.

\textbf{Part (ii).} \textit{Existence.} Define $\bar{f}([x]):=f(x)$ where $[x]$
denotes the $\sim_q$-class of $x$. This is well-defined: if $x\sim_q y$ then
$f(r_1(x,y))=f\circ r_1(x,y)=f\circ r_2(x,y)=f(y)$, so $f(x)=f(y)$.

\textit{$\bar{f}$ is admissibility-preserving.} Let $B\in\mathcal{A}''$. Then
$\bar{f}^{-1}(B)=\{[x]:\bar{f}([x])\in B\}=\{[x]:f(x)\in B\}=\pi_q(f^{-1}(B))$.
Since $f$ is admissibility-preserving, $f^{-1}(B)\in\mathcal{A}$. By definition of
$\mathcal{A}_q$: $B'\in\mathcal{A}_q$ iff $\pi_q^{-1}(B')\in\mathcal{A}$. We have
$\pi_q^{-1}(\bar{f}^{-1}(B))=f^{-1}(B)\in\mathcal{A}$, hence $\bar{f}^{-1}(B)\in
\mathcal{A}_q$.

\textit{$\bar{f}$ is causal-preorder-preserving.} Let $[x]\preceq_{q}[y]$ in
$\Omega/{\sim_q}$: every $A_q\in\mathcal{A}_q$ containing $[y]$ also contains $[x]$.
We must show $\bar{f}([x])\preceq''\bar{f}([y])$ in $\Omega''$. Let $B\in\mathcal{A}''$
with $\bar{f}([y])\in B$, i.e. $f(y)\in B$. Since $f$ is causal-preorder-preserving
and $y\in f^{-1}(B)\in\mathcal{A}$, we need $x\in f^{-1}(B)$. Now
$f^{-1}(B)\in\mathcal{A}$, so $\pi_q(f^{-1}(B))\in\mathcal{A}_q$. We have
$[y]\in\pi_q(f^{-1}(B))$ and by $[x]\preceq_q[y]$ we get $[x]\in\pi_q(f^{-1}(B))$,
i.e. $x\in f^{-1}(B)$, i.e. $f(x)=\bar{f}([x])\in B$.

\textit{The equation $\bar{f}\circ\Col_q=f$ holds} by construction: $\bar{f}(\pi_q(x))
=\bar{f}([x])=f(x)$.

\textit{Uniqueness.} If $\bar{f}'$ also satisfies $\bar{f}'\circ\Col_q=f$, then for
any $[x]\in\Omega/{\sim_q}$: $\bar{f}'([x])=\bar{f}'(\pi_q(x))=f(x)=\bar{f}([x])$.
Since $\pi_q$ is surjective, $\bar{f}'=\bar{f}$.

\textbf{Part (iii).} For any morphism $g:(\Omega'',\mathcal{A}'')\to
(\Omega''',\mathcal{A}''')$, the natural transformation condition
$\overline{g\circ f}=g\circ\bar{f}$ holds because:
$\overline{g\circ f}([x])=g(f(x))=g(\bar{f}([x]))=(g\circ\bar{f})([x])$.
\end{proof}

\paragraph{Summary.} Collapse is not an ad hoc quotient but the unique coequalizing
map satisfying the universal property with respect to causal-preorder-preserving morphisms.
Every categorical framework that admits entropy-decreasing quotients must factor through
$\Col$ when the equivalence relation is causally admissible.

% ═════════════════════════════════════════════════════════════════════════════
\section{Symmetric Monoidal Structure and Universal Property of \texorpdfstring{$\SP$}{SP}}
\label{sec:monoidal}

\subsection{Tensor Product}

\begin{definition}[Tensor of option-spaces]
\label{def:tensor}
\[
  (\Omega_1,\mathcal{A}_1)\tensor(\Omega_2,\mathcal{A}_2)
  := (\Omega_1\times\Omega_2,\;\mathcal{A}_1\tensor\mathcal{A}_2),
\]
where $\mathcal{A}_1\tensor\mathcal{A}_2:=\{U_1\times U_2:U_1\in\mathcal{A}_1,
U_2\in\mathcal{A}_2\}$. The unit object is $\unit=(\{*\},\{\{*\}\})$.

On morphisms: $(e_1\tensor e_2)(x_1,x_2)=(e_1(x_1),e_2(x_2))$.
\end{definition}

\begin{proposition}[Symmetric monoidal structure]
\label{prop:monoidal}
$(\SP,\tensor,\unit)$ is a symmetric monoidal category. By Mac~Lane's
coherence theorem, it may be treated as strict without loss of generality.
\end{proposition}

\begin{proof}
The tensor product is given on objects by Cartesian product and on
morphisms by componentwise application. Associativity and unit
constraints arise from the canonical bijections
\[
(\Omega_1\times\Omega_2)\times\Omega_3
\cong
\Omega_1\times(\Omega_2\times\Omega_3),
\qquad
\Omega\times\{*\}\cong\Omega\cong\{*\}\times\Omega,
\]
which preserve admissibility by functoriality of products.

Symmetry is given by the swap map
\[
\sigma_{\Omega_1,\Omega_2}(x_1,x_2)=(x_2,x_1),
\]
which preserves admissible sets and is involutive, $\sigma^2=\id$.
Entropy is invariant under factor permutation, so symmetry is
entropy-neutral.

Functoriality follows since, for admissibility-preserving morphisms
$e_i$,
\[
(e_1\tensor e_2)^{-1}(U_1\times U_2)
=
e_1^{-1}(U_1)\times e_2^{-1}(U_2)
\in \mathcal{A}_1\tensor\mathcal{A}_2.
\]

The interchange law
\[
(f_2\circ f_1)\tensor(g_2\circ g_1)
=
(f_2\tensor g_2)\circ(f_1\tensor g_1)
\]
holds pointwise:
\[
(x,y)\mapsto
\bigl(f_2(f_1(x)),\,g_2(g_1(y))\bigr).
\]

Since Cartesian product of sets is strictly associative and unital,
the structural isomorphisms may be identified with identities, yielding
a strict symmetric monoidal structure.
\end{proof}

\begin{proposition}[Entropy subadditivity]
\label{prop:subadditive}
\[
\Ent(\Omega_1\tensor\Omega_2)
\le
\Ent(\Omega_1)+\Ent(\Omega_2),
\]
with equality if and only if the joint measure on
$\Omega_1\times\Omega_2$ is a product measure.
\end{proposition}

\begin{proof}
Let $p(x,y)$ be a joint probability distribution with marginals
$p_1(x)$ and $p_2(y)$. Then
\begin{align*}
\Ent(\Omega_1\tensor\Omega_2)
&= -\sum_{x,y} p(x,y)\log p(x,y) \\
&= -\sum_{x,y} p(x,y)\log p_1(x)
   -\sum_{x,y} p(x,y)\log\frac{p(x,y)}{p_1(x)} \\
&= \Ent(\Omega_1) + \Ent(\Omega_2\mid\Omega_1).
\end{align*}
Since conditional entropy satisfies
\[
\Ent(\Omega_2\mid\Omega_1)\le \Ent(\Omega_2),
\]
we obtain the stated inequality. Equality holds precisely when
$p(x,y)=p_1(x)p_2(y)$, i.e., when $\Omega_1$ and $\Omega_2$
are independent.
\end{proof}

\subsection{The Category \texorpdfstring{$\EDSMC$}{EDSMC} and Initiality of \texorpdfstring{$\SP$}{SP}}

\begin{definition}[EDSMC]
\label{def:edsmc}
The category $\EDSMC$ of \emph{entropy-decreasing symmetric monoidal categories}
has:
\begin{itemize}
\item \textit{Objects:} tuples $(\mathcal{C},\tensor_\mathcal{C},\unit_\mathcal{C},
  \Ent_\mathcal{C},P,R,B,C)$ where $(\mathcal{C},\tensor_\mathcal{C},\unit_\mathcal{C})$
  is a strict symmetric monoidal category, $\Ent_\mathcal{C}:\Ob(\mathcal{C})\to\Rnn$
  is monotone along all morphisms, and $P,R,B,C$ are families of morphisms in
  $\mathcal{C}$ satisfying entropy axioms analogous to those of
  $\Pop,\RefOp,\Bind,\Col$ in $\SP$.
\item \textit{Morphisms:} strict symmetric monoidal functors
  $\Phi:\mathcal{C}\to\mathcal{D}$ that preserve $\Ent$ (i.e.
  $\Ent_\mathcal{D}(\Phi(X))=\Ent_\mathcal{C}(X)$ for all $X$) and map generators
  to generators ($\Phi(P_\mathcal{C})\subseteq P_\mathcal{D}$, etc.).
\end{itemize}
\end{definition}

\begin{theorem}[Initiality of $\SP$]
\label{thm:initial}
$\SP$ is the initial object in $\EDSMC$: for every $(\mathcal{C},\ldots)\in\EDSMC$,
there exists a unique morphism $F_\mathcal{C}:\SP\to\mathcal{C}$ in $\EDSMC$.
\end{theorem}

\begin{proof}
\textit{Existence.} We construct $F_\mathcal{C}$ explicitly. On objects: send an
option-space $(\Omega,\mathcal{A})$ to the image of the initial construction in
$\mathcal{C}$. Since $\SP$ is freely generated (Definition~\ref{def:free-gen}), a
morphism out of $\SP$ is determined by its values on generators, subject to the
relations in Definition~\ref{def:free-gen}. Define $F_\mathcal{C}(\Pop_U)=P_U^{(\mathcal{C})}$
(the corresponding $P$-family morphism in $\mathcal{C}$), and similarly for $\RefOp,\Bind,\Col$.
The relations (a)--(d) of Definition~\ref{def:free-gen} hold in $\mathcal{C}$ by hypothesis
(since $\mathcal{C}\in\EDSMC$), so $F_\mathcal{C}$ is well-defined on all morphisms.
$F_\mathcal{C}$ is strict symmetric monoidal because $\SP$'s monoidal structure is
free and $\mathcal{C}$'s monoidal structure satisfies the same strict laws.
Entropy preservation: $\Ent_\mathcal{C}(F_\mathcal{C}(\Omega))=\Ent(\Omega)$ by
construction (the $\Ent$-functional in $\mathcal{C}$ is required to agree with
the abstract entropy assigned to each generator's domain and codomain).

\textit{Uniqueness.} Any morphism $G:\SP\to\mathcal{C}$ in $\EDSMC$ must send
generators to generators and preserve composition. Since $\SP$ is freely generated,
$G$ is completely determined by its values on generators. The generator values are
uniquely forced by the type constraints: $G(\Pop_U)$ must be a $P$-family morphism
at $G(\Omega,\mathcal{A})$, and there is exactly one such morphism for each
elimination $U$ by the hypothesis that $\mathcal{C}$'s generators satisfy the same
axioms as $\SP$'s. Hence $G=F_\mathcal{C}$.
\end{proof}

\begin{remark}
Theorem~\ref{thm:initial} upgrades $\SP$ from ``a constructed category'' to ``the
initial object in the class of entropy-decreasing systems.'' Every such system
receives a canonical map from $\SP$, and this map is unique. The construction of the
realization functor $F:\SP\to\RSVP$ (Section~\ref{sec:functor}) is a special case of
this universal property.
\end{remark}

\paragraph{Summary.} $\SP$ carries a strict symmetric monoidal structure independent
of its sequential composition, with entropy subadditivity under tensor and strict
decrease under composition. It is the initial object in $\EDSMC$: the free
entropy-decreasing rewriting category.

\begin{theorem}[Structural asymmetry]
There exists no functor $G:\RSVP\to\SP$ such that $G\circ F$ is naturally isomorphic to $\mathrm{id}_\SP$. Equivalently, $F$ admits no right-inverse up to natural isomorphism, and in particular $F$ is not part of an equivalence of categories.
\end{theorem}

\begin{proof}[Sketch]
The functor $F$ forgets geometric degrees of freedom by passing from smooth field data to discrete elimination histories. Distinct smooth coupling structures can induce the same discrete history and the same entropy-weight profile, so $F$ is not faithful on the relevant substructures. A natural right-inverse would, by composition, reconstruct the forgotten geometric information from the discrete data in a way stable under morphisms, contradicting the existence of nontrivial diffeomorphism classes and coupling deformations that preserve the induced entropy integrals while varying the smooth structure.
\end{proof}

% ═════════════════════════════════════════════════════════════════════════════
\section{Meld and Sheaf-Theoretic Worldhood}
\label{sec:sheaf}

\subsection{Presheaf of Local Proposals}

Let $(\Omega,\mathcal{A})$ be an option-space and $\mathcal{U}=\{U_i\}_{i\in I}$ a
cover of $\Omega$ by admissible subsets. The \emph{site} $\mathcal{U}$ has the
induced admissibility structure as its topology (finite intersections are admissible).

\begin{definition}[Presheaf of proposals]
\label{def:presheaf}
Define a presheaf $\mathcal{T}:\mathcal{U}^{\mathrm{op}}\to\mathbf{Set}$ by:
\[
  \mathcal{T}(U) = \{\text{admissible event proposals with support in }U\},
\]
with restriction maps $\rho_{UV}:\mathcal{T}(U)\to\mathcal{T}(V)$ for $V\subseteq U$
given by restricting a proposal to the smaller scope.
\end{definition}

A family $\{\delta_i\in\mathcal{T}(U_i)\}_{i\in I}$ is \emph{compatible on overlaps}
if $\rho_{U_i,U_i\cap U_j}(\delta_i)=\rho_{U_j,U_i\cap U_j}(\delta_j)$ for all
$i,j\in I$.

\begin{definition}[Worldhood]
\label{def:worldhood}
$(\Omega,\mathcal{A})$ \emph{exhibits worldhood} with respect to $\mathcal{U}$ if
$\mathcal{T}$ is a sheaf: every compatible family $\{\delta_i\}$ determines a unique
global section $\Delta\in\mathcal{T}(\Omega)$ with $\rho_{\Omega,U_i}(\Delta)=\delta_i$.
\end{definition}

Worldhood is the mathematical content of the claim that locally coherent commitments
can be uniquely assembled into a globally coherent history. Before sufficient
commitment, $\mathcal{T}$ is a presheaf (local sections may fail to glue). Each
Pop, Bind, or Refuse event tightens constraints and moves $\mathcal{T}$ closer to the
sheaf condition.

\subsection{Meld as Policy-Induced Sheafification}

Fix a merge policy $\pi$, which specifies an equivalence relation $\equiv_\pi$ on
proposal data. Let $R_\pi\rightrightarrows\mathcal{T}$ be the corresponding presheaf
equivalence relation. Denote by $q_\pi:\mathcal{T}\to\mathcal{T}/{\equiv_\pi}$ the
pointwise quotient presheaf.

Assume the site is subcanonical (which holds in the admissible-cover setting, since
representable presheaves are sheaves). Let $a:\PSh\to\Sh$ be sheafification, left
adjoint to the inclusion $i:\Sh\hookrightarrow\PSh$.

\begin{definition}[Policy sheafification]
\label{def:policy-sheaf}
Define $a_\pi(\mathcal{T}):=a(\mathcal{T}/{\equiv_\pi})$ and let
$\eta_\pi:\mathcal{T}\to i\,a_\pi(\mathcal{T})$ be the composite of the quotient
$q_\pi$ with the unit $\eta$ of the adjunction.
\end{definition}

\begin{theorem}[Universal property of policy-induced sheafification]
\label{thm:sheafification}
For every sheaf $\mathcal{S}\in\Sh(\mathcal{U})$, composition with $\eta_\pi$ gives
a natural bijection:
\[
  \Hom_{\Sh}(a_\pi(\mathcal{T}),\mathcal{S})
  \;\cong\;
  \{f\in\Hom_{\PSh}(\mathcal{T},i\mathcal{S}):f\text{ is }\pi\text{-invariant}\},
\]
where $\pi$-invariant means $f$ coequalizes $R_\pi\rightrightarrows\mathcal{T}$
(constant on $\equiv_\pi$-classes in each section, stably under restriction).
Equivalently, $a_\pi(\mathcal{T})$ is initial among sheaves receiving a $\pi$-invariant
map from $\mathcal{T}$.
\end{theorem}

\begin{proof}
\textit{Step 1.} By the universal property of coequalizers in $\PSh$: giving a
$\pi$-invariant presheaf morphism $f:\mathcal{T}\to i\mathcal{S}$ is equivalent to
giving a unique presheaf morphism $f':\mathcal{T}/{\equiv_\pi}\to i\mathcal{S}$
with $f'\circ q_\pi=f$.

\textit{Step 2.} By the sheafification adjunction $a\dashv i$:
$\Hom_\Sh(a(\mathcal{P}),\mathcal{S})\cong\Hom_\PSh(\mathcal{P},i\mathcal{S})$,
naturally in $\mathcal{P}$ and $\mathcal{S}$. Apply with $\mathcal{P}=
\mathcal{T}/{\equiv_\pi}$: $\Hom_\Sh(a_\pi(\mathcal{T}),\mathcal{S})
\cong\Hom_\PSh(\mathcal{T}/{\equiv_\pi},i\mathcal{S})$.

\textit{Step 3.} Composing the bijections of Steps 1 and 2:
$\Hom_\Sh(a_\pi(\mathcal{T}),\mathcal{S})\cong$ $\pi$-invariant maps
$\mathcal{T}\to i\mathcal{S}$. The composite unit $\eta_\pi=\eta\circ q_\pi$ is the
universal $\pi$-invariant map.

\textit{Naturality} in $\mathcal{S}$ is inherited from naturality of coequalizer
factorization and of the adjunction bijection.

\textit{Uniqueness of factorization.} Given any sheaf morphism $\bar{f}:
a_\pi(\mathcal{T})\to\mathcal{S}$, the condition $i(\bar{f})\circ\eta_\pi=f$
determines $\bar{f}$ uniquely by the adjunction (the unit $\eta_{\mathcal{T}/\equiv_\pi}$
is initial).
\end{proof}

\begin{definition}[Meld]
\label{def:meld}
Given histories $H_1,H_2$ from a common prefix $\Omega_0$ and a merge policy $\pi$,
the meld $\Meld_\pi(H_1,H_2)$ is the irreversible event in $\SP$ that commits to
the policy-sheafified global section $\Delta^\pi\in a_\pi(\mathcal{T})(\Omega)$ as
the authoritative continuation. The operation decomposes canonically:
\[
  \Meld_\pi \;=\; \mathrm{glue}\circ\Col_\pi,
\]
where $\Col_\pi$ removes the overlap obstructions (the collapse step applying
$\equiv_\pi$) and $\mathrm{glue}$ produces the unique global section
(the sheafification step).
\end{definition}

\begin{remark}
Declaring the global section $\Delta^\pi$ authoritative is not additional mathematical
content within the sheaf category; it is an event in $\SP$ (specifically a $\Col_\pi$
event). The mathematics of Theorem~\ref{thm:sheafification} shows what the
post-meld option-space looks like; the commitment to that option-space is the Spherepop-side
contribution.
\end{remark}

\paragraph{Summary.} Worldhood is the sheaf condition on local event proposals: it is
achieved when commitments are sufficient for local data to glue globally. Meld is the
categorical operation that enforces this by policy-induced sheafification. The theorem
gives it a precise universal property: it is the initial sheaf receiving a $\pi$-invariant
map from the presheaf of proposals.

% ═════════════════════════════════════════════════════════════════════════════
\section{The Category \texorpdfstring{$\RSVP$}{RSVP}}
\label{sec:rsvp}

\subsection{Objects}

\begin{definition}[RSVP state]
\label{def:rsvp-state}
An \emph{RSVP state} is a quadruple $(M,\PhiF,\vF,\SF)$ where $M$ is a smooth,
connected, oriented Riemannian manifold; $\PhiF\in C^\infty(M)$ is the \emph{scalar
coherence potential}; $\vF\in\mathfrak{X}(M)$ is a smooth vector field; and
$\SF\in C^\infty(M,\Rnn)$ is the \emph{entropy density}. The fields may satisfy
the entropy-transport equation $\partial_t\SF=\nabla\cdot(\kappa\nabla\PhiF)$ for
$\kappa>0$, but this is a structural constraint on the dynamics, not a defining
condition on objects.
\end{definition}

\subsection{Morphisms with Slack Witnesses}

\begin{definition}[RSVP morphism]
\label{def:rsvp-morph}
A morphism $(M,\PhiF,\vF,\SF)\to(M',\PhiF',\vF',\SF')$ in $\RSVP$ is a pair
$(\varphi,\eta)$ where $\varphi:M\to M'$ is smooth and $\eta\in C^\infty(M,\Rnn)$
is the \emph{entropy-slack witness} satisfying $\varphi^*\SF'=\SF+\eta$ pointwise.
\end{definition}

\begin{lemma}[Identity morphisms]
\label{lem:rsvp-id}
$(\id_M,0):(M,\PhiF,\vF,\SF)\to(M,\PhiF,\vF,\SF)$ is a morphism in $\RSVP$.
\end{lemma}

\begin{proof}
$\id_M^*\SF=\SF=\SF+0$.
\end{proof}

\begin{lemma}[Closure under composition]
\label{lem:rsvp-compose}
If $(\varphi_1,\eta_1):(M_1,\ldots)\to(M_2,\ldots)$ and $(\varphi_2,\eta_2):
(M_2,\ldots)\to(M_3,\ldots)$ are morphisms in $\RSVP$, then
\[
  (\varphi_2\circ\varphi_1,\;\eta_1+\varphi_1^*\eta_2)
  :(M_1,\ldots)\to(M_3,\ldots)
\]
is a morphism in $\RSVP$.
\end{lemma}

\begin{proof}
$(\varphi_2\circ\varphi_1)^*\SF_3=\varphi_1^*(\varphi_2^*\SF_3)=\varphi_1^*(\SF_2+\eta_2)
=\varphi_1^*\SF_2+\varphi_1^*\eta_2=(\SF_1+\eta_1)+\varphi_1^*\eta_2
=\SF_1+(\eta_1+\varphi_1^*\eta_2)$.
Non-negativity: $\eta_1\ge 0$, $\eta_2\ge 0$, pullback preserves non-negativity.
\end{proof}

\begin{lemma}[Associativity of composition]
\label{lem:rsvp-assoc}
Composition in $\RSVP$ is associative.
\end{lemma}

\begin{proof}
$(\varphi_3\circ(\varphi_2\circ\varphi_1),\;\eta_1+\varphi_1^*\eta_2+\varphi_1^*\varphi_2^*
\eta_3)=((\varphi_3\circ\varphi_2)\circ\varphi_1,\;\eta_1+\varphi_1^*(\eta_2+\varphi_2^*
\eta_3))$ since smooth map composition is associative and $\varphi_1^*$ is additive.
\end{proof}

\begin{proposition}[$\RSVP$ is a category]
\label{prop:rsvp-cat}
$\RSVP$ with composition defined by Lemma~\ref{lem:rsvp-compose} is a well-defined
category.
\end{proposition}

\begin{proof}
Identity morphisms exist by Lemma~\ref{lem:rsvp-id}. Composition is closed
(Lemma~\ref{lem:rsvp-compose}) and associative (Lemma~\ref{lem:rsvp-assoc}).
Left and right unit laws: $(\id,0)\circ(\varphi,\eta)=(\varphi,\eta+\varphi^*0)
=(\varphi,\eta)$; $(\varphi,\eta)\circ(\id,0)=(\varphi,0+\varphi^*|_{\id}\eta)
=(\varphi,\eta)$.
\end{proof}

\paragraph{Summary.} $\RSVP$ is a well-defined category with objects being smooth
entropy-carrying manifolds and morphisms being smooth maps with explicit non-negative
entropy-slack witnesses. The slack data resolve the earlier ambiguity about
``preservation up to boundary terms'' by making the entropy production at each
morphism explicit.

% ═════════════════════════════════════════════════════════════════════════════
\section{The Realization Functor \texorpdfstring{$F:\SP\to\RSVP$}{F}}
\label{sec:functor}

\subsection{On Objects}

\begin{definition}[Geometric realization on objects]
\label{def:F-objects}
For an option-space $(\Omega,\mathcal{A})$, define
\[
  F(\Omega,\mathcal{A}) := (\Delta(\Omega),\,\PhiF_\Omega,\,\vF_\Omega,\,\SF_\Omega),
\]
where $\Delta(\Omega)$ is the \emph{probability simplex} $\{p\in\Rnn^\Omega:\sum_x p(x)=1\}$;
$\SF_\Omega(p)=-\sum_{x\in\Omega}p(x)\log p(x)$ is the Shannon entropy density at
the interior of $\Delta(\Omega)$; $\PhiF_\Omega$ and $\vF_\Omega$ are smooth
interpolants of the admissibility structure, defined by a fixed smooth embedding
of $\mathcal{A}$ into $C^\infty(\Delta(\Omega))$.
\end{definition}

\subsection{On Generating Morphisms}

\paragraph{$F$ on Pop.}
Let $\Pop_U:(\Omega,\mathcal{A})\to(\Omega\setminus U,\cdot)$.

\begin{definition}[Pop realization]
\label{def:F-pop}
Define $F(\Pop_U)=(\iota_U,\eta_U)$ where $\iota_U:\Delta(\Omega\setminus U)
\hookrightarrow\Delta(\Omega)$ is the face inclusion (restriction to the face
corresponding to $\Omega\setminus U$) and $\eta_U\in C^\infty(\Delta(\Omega\setminus
U),\Rnn)$ is a bump function supported near $\partial\Delta(U)$, normalized so that
$\iota_U^*\SF_\Omega=\SF_{\Omega\setminus U}+\eta_U$.
\end{definition}

The bump $\eta_U$ records the entropy difference between the full simplex and the
restricted face. It is the \emph{boundary-sharpening slack}: the entropy production
localized at the boundary of the eliminated region.

\paragraph{$F$ on Refuse.}
$F(\RefOp_U)=F(\Pop_U)$ as a smooth map, with the same slack $\eta_U$. The accounting
tag is preserved by $\mathcal{B}:\SP\to\mathbf{Set}$ and does not affect the field data.

\paragraph{$F$ on Bind.}
Let $\Bind_{ij}:(\Omega,\mathcal{A})\to(\Omega,\mathcal{A}')$.

\begin{definition}[Bind realization]
\label{def:F-bind}
Define $F(\Bind_{ij})=(\id_{\Delta(\Omega)},0)$ with the modification
$\vF_\Omega\mapsto\vF_\Omega+\delta\vF_{ij}$, where $\delta\vF_{ij}$ is a smooth
directed vector field increment coupling the regions $\Delta(\{i\})$ and
$\Delta(\{j\})$ in $\Delta(\Omega)$.
\end{definition}

Entropy slack is zero ($\eta=0$) since $\Bind$ preserves entropy globally. The
coupling $\delta\vF_{ij}$ breaks the rotational symmetry between the $i$- and
$j$-directions of the simplex, realizing the admissibility reduction
$\mathcal{A}'\subsetneq\mathcal{A}$ as a directional bias in the flow.

\paragraph{$F$ on Collapse.}
Let $\Col_\sim:(\Omega,\mathcal{A})\to(\Omega/{\sim},\mathcal{A}_\sim)$.

\begin{definition}[Collapse realization]
\label{def:F-col}
Define $F(\Col_\sim)=(\varphi_\sim,\eta_\sim)$ where $\varphi_\sim:\Delta(\Omega)
\to\Delta(\Omega/{\sim})$ is the coarse-graining map (marginalizing over equivalence
classes: $\varphi_\sim(p)([x]):=\sum_{y\sim x}p(y)$) and $\eta_\sim\in C^\infty
(\Delta(\Omega),\Rnn)$ is the fine-to-coarse entropy difference $\eta_\sim(p)=
\SF_\Omega(p)-\varphi_\sim^*\SF_{\Omega/\sim}(p)\ge 0$ (non-negative by Jensen's
inequality applied to the logarithm).
\end{definition}

This coarse-graining is a Wilsonian renormalization step: fine-grained distinctions
within each equivalence class of $\sim$ are integrated out. The slack $\eta_\sim$
records the entropy of the intra-class distribution.

\subsection{Functoriality}

\begin{lemma}[$F$ preserves identities]
\label{lem:F-id}
$F(\id_{(\Omega,\mathcal{A})})=(\id_{\Delta(\Omega)},0)=\id_{F(\Omega,\mathcal{A})}$.
\end{lemma}

\begin{proof}
The identity history is the empty sequence. Its realization is $(\id,0)$ by
Definition~\ref{def:F-objects}: the simplex is unchanged, no slack is produced.
\end{proof}

\begin{lemma}[$F$ preserves generator composition]
\label{lem:F-comp}
For composable generators $e_1$ and $e_2$:
$F(e_2\circ e_1)=F(e_2)\circ F(e_1)$ with slacks composed additively per
Lemma~\ref{lem:rsvp-compose}.
\end{lemma}

\begin{proof}
We check each combination. The key case is $\Pop_V\circ\Pop_U$ (elimination of $U$
then $V\subseteq\Omega\setminus U$): $F(\Pop_V)\circ F(\Pop_U)=(\iota_V,\eta_V)\circ
(\iota_U,\eta_U)=(\iota_V\circ\iota_U,\eta_U+\iota_U^*\eta_V)$. On the other side,
$F(\Pop_{U\cup V})=(\iota_{U\cup V},\eta_{U\cup V})$ where $\iota_{U\cup V}=\iota_V
\circ\iota_U$ (composition of face inclusions) and $\eta_{U\cup V}(p)=\eta_U(p)+
\eta_V(\iota_U(p))=\eta_U+\iota_U^*\eta_V$. The cases $\Bind\circ\Pop$, $\Col\circ\Pop$,
and combinations thereof follow analogously from the locality of each perturbation
and the functoriality of pullback.
\end{proof}

\begin{lemma}[Slack composition is associative]
\label{lem:slack-assoc}
For three composable morphisms $(\varphi_1,\eta_1),(\varphi_2,\eta_2),(\varphi_3,\eta_3)$:
$(\eta_1+\varphi_1^*\eta_2)+\varphi_1^*\varphi_2^*\eta_3=
\eta_1+\varphi_1^*(\eta_2+\varphi_2^*\eta_3)$.
\end{lemma}

\begin{proof}
By distributivity of $\varphi_1^*$ over addition and functoriality:
$\varphi_1^*(\eta_2+\varphi_2^*\eta_3)=\varphi_1^*\eta_2+\varphi_1^*\varphi_2^*\eta_3$.
\end{proof}

\begin{lemma}[Tensor compatibility]
\label{lem:F-tensor}
$F((\Omega_1,\mathcal{A}_1)\tensor(\Omega_2,\mathcal{A}_2))\cong
F(\Omega_1,\mathcal{A}_1)\tensor_\RSVP F(\Omega_2,\mathcal{A}_2)$,
where the tensor in $\RSVP$ is the product of smooth manifolds with product fields.
\end{lemma}

\begin{proof}
$\Delta(\Omega_1\times\Omega_2)\cong\Delta(\Omega_1)\times\Delta(\Omega_2)$ as smooth
manifolds (a probability measure on a product is a joint measure; the simplex over
the product is the product of simplices). Shannon entropy of the product is the sum
of marginal entropies in the independent case, matching subadditivity.
\end{proof}

\begin{theorem}[Functoriality of $F$]
\label{thm:functor}
$F:\SP\to\RSVP$ is a well-defined symmetric monoidal functor. Specifically:
\begin{enumerate}[label=(\alph*)]
\item $F$ preserves identities (Lemma~\ref{lem:F-id});
\item $F$ preserves composition, with slacks composed additively (Lemma~\ref{lem:F-comp});
\item Slack composition is associative (Lemma~\ref{lem:slack-assoc});
\item $F$ is symmetric monoidal (Lemma~\ref{lem:F-tensor});
\item Entropy compatibility: $\int_{\Delta(\Omega)}\SF_\Omega\,d\mathrm{vol}=\Ent(\Omega)$
  up to normalization, so entropy decrease in $\SP$ corresponds to net positive slack
  in $\RSVP$.
\end{enumerate}
\end{theorem}

\begin{proof}
Parts (a)--(d) follow from the four lemmas above. For (e): by construction,
$\SF_\Omega(p)=-\sum_x p(x)\log p(x)$. Integrating over the uniform measure on
$\Delta(\Omega)$ yields $\Ent(\Omega)$ (up to the normalization factor of the simplex
volume, which is absorbed into the definition of the volume form). Entropy decrease
$\Ent(\Omega')<\Ent(\Omega)$ after $\Pop_U$ corresponds to $\eta_U>0$ being the
positive slack that records the entropy difference between $\SF_\Omega$ and
$\iota_U^*\SF_\Omega$.
\end{proof}

\paragraph{Commutative diagram.}
\[
\begin{tikzcd}[column sep=4.5em, row sep=3em]
  (\Omega,\mathcal{A})
  \arrow[r,"e"]
  \arrow[d,"F"']
  &
  (\Omega',\mathcal{A}')
  \arrow[d,"F"]
  \\
  (\Delta(\Omega),\PhiF,\vF,\SF)
  \arrow[r,"{(\varphi,\eta)}"']
  &
  (\Delta(\Omega'),\PhiF',\vF',\SF')
\end{tikzcd}
\]


\section{Initiality of \texorpdfstring{$\SP$}{SP}: Mac Lane Presentation}
\label{sec:initiality-maclane}

We refactor Theorem~\ref{thm:initial} (initiality of $\SP$ in $\EDSMC$) in the
style of Mac Lane's treatment of free monoidal categories \citep{maclane1998}.
The benefit is a clean separation between the \emph{term model} (syntax),
the \emph{congruence} (axioms), and the \emph{evaluation functor} (semantics),
making uniqueness immediate without appeal to ``exactly one generator per type.''

\subsection{The Term Model \texorpdfstring{$\mathcal{T}$}{T}}

\begin{definition}[Term model]
\label{def:term-model}
Let $\mathsf{Ob}$ be the class of all option-spaces. Define the \emph{term model}
$\mathcal{T}$ as follows.
\begin{itemize}
  \item \emph{Objects.} Formal tensor-products of elements of $\mathsf{Ob}$ under
    a binary operation $\otimes$ and a unit symbol $\mathsf{I}$: the free commutative
    monoid generated by $\mathsf{Ob}$.
  \item \emph{Morphisms.} Formal expressions built from:
    \begin{itemize}
      \item identity symbols $\mathsf{id}_X$ for each object $X$;
      \item generator symbols $\mathsf{Pop}_U$, $\mathsf{Ref}_U$ for each
        nonempty $U\subsetneq\Omega$;
      \item generator symbols $\mathsf{Bind}_{ij}$ for each valid dependency pair;
      \item generator symbols $\mathsf{Col}_\sim$ for each causally admissible $\sim$;
    \end{itemize}
    closed under composition $\circ$ and tensor $\otimes$ with sources and targets
    determined by the set-theoretic typing already given in
    Definitions~\ref{def:pop}--\ref{def:col}.
\end{itemize}
$\mathcal{T}$ is a strict symmetric monoidal category (the free such category on
these generators and types, before imposing any axioms beyond typing).
\end{definition}

\subsection{The Congruence \texorpdfstring{$\equiv$}{≡}}

\begin{definition}[Presentation congruence]
\label{def:congruence}
Let $\equiv$ be the smallest congruence on morphisms of $\mathcal{T}$ containing:
\begin{enumerate}[label=(\alph*)]
  \item the strict symmetric monoidal axioms (associativity, unit, symmetry,
    interchange law);
  \item \emph{Pop-idempotence:}
    $\mathsf{Pop}_U\circ\mathsf{Pop}_V\equiv\mathsf{Pop}_{U\cup V}$
    when $U\cap(\Omega\setminus V)=\emptyset$;
  \item \emph{Independent-bind commutativity:}
    $\mathsf{Bind}_{ij}\circ\mathsf{Bind}_{kl}\equiv
    \mathsf{Bind}_{kl}\circ\mathsf{Bind}_{ij}$ when $\{i,j\}\cap\{k,l\}=\emptyset$;
  \item \emph{Entropy monotonicity:} any morphism $f$ with $\Ent(\mathrm{cod}\,f)>
    \Ent(\mathrm{dom}\,f)$ is identified with the empty class (i.e., such expressions
    are inadmissible and removed from $\mathcal{T}$).
\end{enumerate}
\end{definition}

\begin{proposition}
$\SP\cong\mathcal{T}/{\equiv}$ as strict symmetric monoidal categories.
\end{proposition}

\begin{proof}
The quotient $\mathcal{T}/{\equiv}$ has the same objects and generating morphisms
as the category defined in Section~\ref{sec:sp}, modulo exactly the relations of
Definition~\ref{def:free-gen}. The identification is the identity on generators.
\end{proof}

\subsection{Evaluation Functors and Initiality}

\begin{definition}[EDSMC (revised)]
\label{def:edsmc-revised}
An object of $\EDSMC$ is a strict symmetric monoidal category $\mathcal{C}$ together
with distinguished families of morphisms
$\mathbf{P},\mathbf{R},\mathbf{B},\mathbf{C}$ interpreting the four generator
symbols, and a monotone functional $\Ent_\mathcal{C}:\Ob(\mathcal{C})\to\Rnn$,
such that the relations (a)--(d) of Definition~\ref{def:congruence} hold
(with generator families substituted for generator symbols). A morphism in
$\EDSMC$ is a strict symmetric monoidal functor that maps each generator family
to the corresponding generator family and satisfies
$\Ent_\mathcal{C}(X)=\Ent_{\mathcal{D}}(\Phi(X))$ for all $X$.
\end{definition}

\begin{theorem}[Initiality, Mac Lane presentation]
\label{thm:initial-maclane}
$\SP=\mathcal{T}/{\equiv}$ is the initial object in $\EDSMC$: for every
$\mathcal{C}\in\EDSMC$ there is a unique morphism $\Phi_\mathcal{C}:\SP\to\mathcal{C}$
in $\EDSMC$.
\end{theorem}

\begin{proof}
\emph{Existence.} Define $\Phi_\mathcal{C}$ on generators by
\[
  \Phi_\mathcal{C}(\mathsf{Pop}_U):=P^{(\mathcal{C})}_U,\quad
  \Phi_\mathcal{C}(\mathsf{Ref}_U):=R^{(\mathcal{C})}_U,\quad
  \Phi_\mathcal{C}(\mathsf{Bind}_{ij}):=B^{(\mathcal{C})}_{ij},\quad
  \Phi_\mathcal{C}(\mathsf{Col}_\sim):=C^{(\mathcal{C})}_\sim.
\]
Extend by strict monoidality. Well-definedness on $\mathcal{T}/{\equiv}$: the
relations (a)--(d) hold in $\mathcal{C}$ by hypothesis, so $\Phi_\mathcal{C}$
respects $\equiv$.

\emph{Uniqueness.} Any morphism $\Psi:\SP\to\mathcal{C}$ in $\EDSMC$ must send
each generator family to the corresponding family in $\mathcal{C}$ (by the
definition of $\EDSMC$-morphisms) and be strictly monoidal. Since $\SP$ is
generated by the four families under composition and tensor, $\Psi$ is uniquely
determined by these values. Hence $\Psi=\Phi_\mathcal{C}$.
\end{proof}

\begin{remark}
Theorem~\ref{thm:initial-maclane} supersedes Theorem~\ref{thm:initial} with a
cleaner proof structure. Uniqueness no longer requires the ``exactly one generator
per elimination'' argument; it follows immediately from the fact that
$\EDSMC$-morphisms are required to preserve generator families, and $\SP$ is freely
generated by those families. The realization functor $F:\SP\to\RSVP$
(Section~\ref{sec:functor}) is the special case $\mathcal{C}=\RSVP$ with
generator families $\mathbf{P}=\{(\iota_U,\eta_U)\}$, $\mathbf{B}=\{(\id,\delta v_{ij})\}$,
$\mathbf{C}=\{(\varphi_\sim,\eta_\sim)\}$.
\end{remark}

\paragraph{Structural asymmetry theorem.}

\begin{theorem}[Structural asymmetry]
\label{thm:asymmetry}
$\SP$ and $\RSVP$ are related by $F$ but are not equivalent categories. The asymmetry
is:
\begin{enumerate}[label=(\roman*)]
\item $\SP$ is discrete and combinatorial; $\RSVP$ is smooth and differential.
\item $\SP$ \emph{accumulates constraint}: each morphism adds an irreversible record;
  no morphism in $\SP$ increases entropy.
\item $\RSVP$ \emph{redistributes coherence}: the entropy-transport equation
  $\partial_t\SF=\nabla\cdot(\kappa\nabla\PhiF)$ diffuses entropy across $M$; local
  concentrations propagate globally.
\item $F$ is not an equivalence of categories: there exist morphisms in $\RSVP$ that
  are not in the image of $F$ (smooth maps that do not correspond to any finite
  composition of the four generators).
\end{enumerate}
\end{theorem}

\begin{proof}
(i) $\SP$ has discrete objects (finite option-spaces); $\RSVP$ has smooth manifolds.
(ii) By Axiom~\ref{ax:entropy}, all morphisms in $\SP$ are entropy-monotone non-increasing.
(iii) The PDE $\partial_t\SF=\nabla\cdot(\kappa\nabla\PhiF)$ is parabolic; it
redistributes entropy according to the gradient of $\PhiF$, potentially increasing
$\SF$ locally (where $\nabla\cdot(\kappa\nabla\PhiF)>0$) while decreasing it elsewhere.
This is not monotone. (iv) The image of $F$ consists of restriction maps, vector field
additions, and coarse-graining maps of the prescribed forms. A generic smooth map
$\varphi:M\to M'$ need not have this structure.
\end{proof}

% ═════════════════════════════════════════════════════════════════════════════
\section{Worked Example: A Four-Event History}
\label{sec:worked}

We trace a complete history through $\SP$ and $\RSVP$.

\subsection{Setup}

Let $\Omega_0=\{a,b,c,d\}$ with uniform probability and trivial admissibility
$\mathcal{A}_0=2^{\Omega_0}$. Initial entropy: $\Ent(\Omega_0)=\log 4=2\log 2$.
Initial optionality: $\Opt(\Omega_0)=2\log 2$. Under $F$: the probability simplex
$\Delta(\Omega_0)$ is a regular 3-simplex with vertices $e_a,e_b,e_c,e_d$.
Initial entropy density $\SF(p)=-\sum_x p(x)\log p(x)$ achieves maximum $\log 4$
at the centroid.

\subsection{Event 1: Pop}

$e_1=\Pop_{\{c,d\}}:(\Omega_0,\mathcal{A}_0)\to(\Omega_1,\mathcal{A}_1)$
where $\Omega_1=\{a,b\}$ and $\mathcal{A}_1=\{\{a,b\},\{a\},\{b\},\emptyset\}$.

\paragraph{In $\SP$.}
$\Ent(\Omega_1)=\log 2$. Entropy decrease: $\Delta\Ent=\log 4-\log 2=\log 2$.
Optionality: $\Opt(\Omega_1)=\log 2$. Action increment: $L_1=\log 2$.

\paragraph{In $\RSVP$.}
$F(e_1)=(\iota_{\{c,d\}},\eta_{\{c,d\}})$. The simplex $\Delta(\Omega_1)$ is the edge
$[e_a,e_b]$ of the original tetrahedron. The slack function $\eta_{\{c,d\}}$ is a
smooth bump on $[e_a,e_b]$ that is positive near the boundary points $e_a$ and $e_b$
and equals $\SF_{\Omega_0}(\iota(p))-\SF_{\Omega_1}(p)$ pointwise, which is the
entropy of the eliminated mass on $\{c,d\}$. Geometrically: the coherence potential
$\PhiF$ sharpens at the boundary of the eliminated region; $\SF$ decreases from
its maximum $\log 4$ toward $\log 2$.

\subsection{Event 2: Bind}

$e_2=\Bind_{ab}:(\Omega_1,\mathcal{A}_1)\to(\Omega_2,\mathcal{A}_2)$
with $\mathcal{A}_2=\{\{a,b\},\{a\},\emptyset\}$ (every admissible future containing
$b$ must contain $a$).

\paragraph{In $\SP$.}
$\Ent(\Omega_2)=\log 2=\Ent(\Omega_1)$: entropy preserved. Optionality unchanged:
$\Opt(\Omega_2)=\log 2$. Action increment: $L_2=0$. But the causal preorder is
now $a\preceq b$ (causally, $a$ is necessary for $b$). The symmetry group of
$(\Omega_2,\mathcal{A}_2)$ has been reduced: the two elements $a,b$ are no longer
equivalent under admissibility-preserving automorphisms.

\paragraph{In $\RSVP$.}
$F(e_2)=(\id_{[e_a,e_b]},0)$ with $\vF\mapsto\vF+\delta\vF_{ab}$, where $\delta\vF_{ab}$
is a smooth vector field on $[e_a,e_b]$ directed from $e_a$ toward $e_b$. This
represents the directed coupling: probability flow is biased toward $b$ (the dependent
element) conditional on $a$ (the prerequisite). The entropy density $\SF$ is unchanged;
the coherence structure $\vF$ is modified.

\subsection{Event 3: Renaming (not a Collapse)}

We consider $e_3=\Col_{a\sim b}$ under the equivalence $a\sim b$. Causal admissibility fails:
under the causal preorder induced by Event~2, we have $\downset{a}=\{a\}$ while $\downset{b}=\{a,b\}$, hence $\downset{a}\neq\downset{b}$ and Definition~\ref{def:causal-admissible} forbids this collapse.

\begin{remark}
This is a genuine failure mode: collapsing after binding can violate causal admissibility, because binding changes causal pasts. The allowed repairs within the current generator set are to apply the intended collapse before the binding step, or to choose a different causally admissible collapse policy. If one wishes to model an explicit ``authority relabeling'' that introduces fresh names, that operation is not a quotient and should be added as a separate primitive (or treated as an isomorphism) rather than being folded into $\Col$.
\end{remark}

Accordingly, we take $e_3$ to be a harmless relabeling isomorphism $\rho:\Omega_2\to\Omega_3$ with $\Omega_2=\{a,b\}$ and $\Omega_3=\{a,[b]\}$, defined by $\rho(a)=a$ and $\rho(b)=[b]$, and with admissibility transported along $\rho$ so that $\mathcal{A}_3=\{\rho(A):A\in\mathcal{A}_2\}$.

\paragraph{In $\SP$.}
Entropy and optionality are unchanged: $\Ent(\Omega_3)=\Ent(\Omega_2)$ and $L_3=0$. The step exists only to make the later notation explicit; it performs no quotienting and should not be confused with $\Col$.

\paragraph{In $\RSVP$.}
$F(\rho)$ is the identity on the underlying simplex up to the canonical identification induced by the relabeling of vertices, and the slack remains zero.

\subsection{Event 4: Pop to Terminal}

$e_4=\Pop_{\{[b]\}}:(\Omega_3,\mathcal{A}_3)\to(\Omega_4,\mathcal{A}_4)$
where $\Omega_4=\{a\}$.

\paragraph{In $\SP$.}
$\Ent(\Omega_4)=0$. Entropy decrease: $\Delta\Ent=\log 2$. Action increment: $L_4=\log 2$.
The world is now fully committed: a single future $a$ remains. Total action: $\mathcal{S}[\gamma]
=L_1+L_2+L_3+L_4=\log 2+0+0+\log 2=\log 4=\Ent(\Omega_0)$.

\paragraph{In $\RSVP$.}
$F(e_4)=(\iota_{\{[b]\}},\eta_{\{[b]\}})$ where $\iota$ maps the vertex $e_a$ to
the edge endpoint, and $\eta$ records the final entropy production $\log 2$. The
coherence potential $\PhiF$ concentrates at the vertex $e_a$; entropy density
$\SF\to 0$ at the vertex. This is the \emph{coherence attractor}: the trajectory has
converged to a point.

\subsection{Summary of the History}

\begin{center}
\renewcommand{\arraystretch}{1.5}
\begin{tabular}{lllll}
\toprule
\textbf{Step} & \textbf{Event} & $\Ent$ & $L_t$ & \textbf{RSVP image} \\
\midrule
0 & --- & $\log 4$ & --- & Centroid of tetrahedron \\
1 & $\Pop_{\{c,d\}}$ & $\log 2$ & $\log 2$ & Restriction to edge $[e_a,e_b]$ \\
2 & $\Bind_{ab}$ & $\log 2$ & $0$ & Directed flow $a\to b$ \\
3 & $\Col_{\sim_{\mathrm{triv}}}$ & $\log 2$ & $0$ & Identity (trivial collapse) \\
4 & $\Pop_{\{[b]\}}$ & $0$ & $\log 2$ & Convergence to vertex $e_a$ \\
\midrule
\multicolumn{3}{l}{Total action} & $\log 4$ & \\
\bottomrule
\end{tabular}
\end{center}

The total action equals the initial entropy, as expected: a history that fully commits
a world of entropy $\log 4$ must pay action $\log 4$.

% ═════════════════════════════════════════════════════════════════════════════
\section{Discrete Mechanics of Commitment}
\label{sec:mechanics}

\subsection{Lagrangian Formulation}

Each event $e_t:X_{t-1}\to X_t$ induces a local cost
$L_t=\Delta\Opt_t+\lambda\,\Delta C_t$ where
$\Delta\Opt_t=\Opt(X_{t-1})-\Opt(X_t)$ and $\Delta C_t$ records auxiliary
accounting costs. The action of a history is
$\mathcal{S}[\gamma]=\sum_{t=1}^n L_t.$

\begin{proposition}
Every history with at least one $\Pop$ or $\RefOp$ event has $\mathcal{S}[\gamma]>0$.
\end{proposition}

\begin{proof}
$\Pop$ and $\RefOp$ strictly reduce $\Opt$, contributing $\Delta\Opt_t>0$. Since all
terms are non-negative, $\mathcal{S}[\gamma]>0$.
\end{proof}

\subsection{Collapse Reduces Action}

If $\Col$ identifies $k$ elements, action reduces by $\log k$ in the Shannon
realization: $\mathcal{S}\mapsto\mathcal{S}-\log k$.

This matches the RSVP side: the slack produced by $F(\Col_\sim)$ equals the
fine-structure entropy integrated out, which is $\log k$ for a uniform distribution
over $k$ identified elements.

\subsection{Hamiltonian and Remaining Freedom}

Define commitment $\pi_t=-\partial L_t/\partial\Opt$ and Hamiltonian
$H_t=\pi_t\Opt(X_t)-L_t$. Then $H_t$ decays monotonically toward zero as commitment
accumulates, measuring the remaining flexibility of the system. A fully committed world
has $\Opt=0$ and $H=0$.


\section{Continuum Limit of Commitment} \label{sec:continuum}

Let $\Delta t \to 0$ and assume optionality is smooth in $t$. Define
\[
\Omega(t+\Delta t) - \Omega(t) = -\alpha \int_{\partial M_t} \nabla\Phi \cdot \mathbf{n}\, dA \,\Delta t.
\]
Then
\[
\frac{d\Omega}{dt} = -\alpha \int_{\partial M_t} \nabla\Phi \cdot \mathbf{n}\, dA.
\]
Using the divergence theorem:
\[
\frac{d\Omega}{dt} = -\alpha \int_{M_t} \Delta \Phi\, dV.
\]
Comparing with entropy transport ($\partial_t S = \nabla\cdot(\kappa\nabla\Phi)$, $\alpha = \kappa$):
optionality decrease is the integral of entropy flux across the boundary---the discrete pop is the boundary flux of coherence.

% ═════════════════════════════════════════════════════════════════════════════
\section{Information-Theoretic Interpretation}
\label{sec:info}

Each generator corresponds to a standard information-theoretic operation:
\begin{align*}
\Pop &:\quad\Ent_{t+1}=\Ent_t-\Delta,\quad\Delta>0 \\
\Bind &:\quad I(X;Y)\uparrow\text{ (mutual information coupling increases)} \\
\Col &:\quad\Kc(\Omega/{\sim})\le\Kc(\Omega)\text{ (Kolmogorov complexity decreases)} \\
F &:\quad\partial_t\SF=\nabla\cdot(\kappa\nabla\PhiF)\text{ (entropy-transport PDE)}
\end{align*}

The second relation captures why $\Bind$ is entropy-neutral at the distribution level
but produces information-theoretic work: it increases the mutual information between
the bound elements without changing marginal entropies, encoding a dependency that was
not there before.

% ═════════════════════════════════════════════════════════════════════════════
\section{Normalization}
\label{sec:norm}

\begin{theorem}[Strong normalization, deterministic fragment]
\label{thm:sn}
Every history in $\SP$ without Choice events, and in which every $\Bind$ step strictly refines admissibility and every $\Col$ step is nontrivial, terminates.
\end{theorem}

\begin{proof}
Define the termination measure $\mu(\Omega,\mathcal{A})=|\Omega|+|\mathcal{A}|$. Each $\Pop$ strictly decreases $|\Omega|$. Each $\Bind$ strictly decreases $|\mathcal{A}|$ by hypothesis (strict refinement). Each nontrivial $\Col$ strictly decreases $|\Omega|$. Hence $\mu$ strictly decreases at each step, and since $\mu$ is a natural number bounded below, the history must be finite.
\end{proof}

\begin{remark}[Probabilistic normalization]
In the presence of Choice events, almost-sure normalization is conjectured: every
well-typed closed probabilistic term terminates with probability 1. The formal proof
requires a probabilistic termination measure satisfying the supermartingale condition
$\mathbb{E}[\mu_{t+1}]\le\mu_t-\delta$ for some $\delta>0$. This is plausible under
bounded entropy conditions but remains open.
\end{remark}

\section{Structured Irreversibility}
\label{sec:structured}

The defining feature of $\SP$ is not entropy monotonicity alone, but the
structural absence of inverses together with resource descent. This
section isolates the precise categorical content of this condition.

\begin{definition}[Entropy preorder]
Let $(\SP,\Ent)$ be as defined previously. Define a preorder
$\preceq$ on $\Ob(\SP)$ by
\[
  X \preceq Y \iff \exists f : Y \to X.
\]
\end{definition}

\begin{proposition}
If $X \preceq Y$, then $\Ent(X) \le \Ent(Y)$.
\end{proposition}

\begin{proof}
Immediate from entropy monotonicity along morphisms.
\end{proof}

Thus entropy induces a monotone map from the preorder $(\Ob(\SP),\preceq)$
into $(\mathbb{R}_{\ge 0},\le)$.

\begin{definition}[Structured irreversibility]
A symmetric monoidal category $(\mathcal{C},\otimes,I)$ equipped with
$R : \Ob(\mathcal{C}) \to \mathbb{R}_{\ge 0}$ is
\emph{structurally irreversible} if:
\begin{enumerate}
  \item For all non-identity $f$, there exists no $g$ with
        $g\circ f=\id$ or $f\circ g=\id$.
  \item $R$ is monotone along morphisms.
  \item $R(X)=0$ implies $X$ is minimal in the preorder.
\end{enumerate}
\end{definition}

\begin{theorem}
$\SP$ is structurally irreversible.
\end{theorem}

\begin{proof}
(1) is Axiom~\ref{ax:irrev}.
(2) is Axiom~\ref{ax:entropy}.
(3) follows from Definition~\ref{def:entropy}(i).
\end{proof}

Hence irreversibility in $\SP$ is not an emergent phenomenon but a
categorical constraint. The resource functional does not merely measure
behavior; it defines directionality.

\section{Irreversibility Across Scales}
\label{sec:scale}

Let $F : \SP \to \RSVP$ be the symmetric monoidal functor constructed
in Section~\ref{sec:functor}. This section formalizes how structural
irreversibility is preserved under $F$.

\begin{proposition}[Entropy preservation under $F$]
For any $X \in \Ob(\SP)$,
\[
  \Ent(X) = \int_{\Delta(X)} S_X \, d\mu
\]
up to normalization.
\end{proposition}

\begin{proof}
By construction of $F(X) = (\Delta(X),S_X)$ with $S_X$ the Shannon entropy density.
Integration over the simplex yields the discrete entropy.
\end{proof}

\begin{proposition}[Slack monotonicity]
For any morphism $f$ in $\SP$,
\[
  F(f) = (\varphi,\eta)
\]
satisfies $\eta \ge 0$ and
\[
  \int \eta \, d\mu = \Ent(X) - \Ent(Y).
\]
\end{proposition}

\begin{proof}
By construction of slack witnesses in Section~\ref{sec:functor}.
\end{proof}

Thus entropy descent in $\SP$ corresponds to non-negative entropy
production in $\RSVP$.

\begin{theorem}[Scale invariance of irreversibility]
If $f$ is non-invertible in $\SP$, then $F(f)$ is non-invertible in
$\RSVP$.
\end{theorem}

\begin{proof}
Suppose $(\varphi,\eta)$ had an inverse. Then slack would have to vanish
and $\varphi$ be a diffeomorphism. This contradicts strict entropy
descent for $\Pop$ and $\RefOp$.
\end{proof}

Therefore, irreversibility is preserved functorially across discrete
and continuous scales.

\section{Persistence Without Substance}
\label{sec:persistence}

Identity in $\SP$ is defined entirely by morphism chains.

\begin{definition}[Log equivalence]
Two histories $\gamma,\gamma'$ are equivalent under policy $q$ if
\[
  q(\gamma(\Omega_0)) = q(\gamma'(\Omega_0)).
\]
\end{definition}

\begin{proposition}
Objects of $\SP$ are equivalence classes of histories modulo $q$.
\end{proposition}

\begin{proof}
By Axiom~\ref{ax:auth} and Axiom~\ref{ax:colid}.
\end{proof}

Hence object identity is induced, not primitive.

Under $F$, identity corresponds to simplex face inclusion:

\begin{proposition}
If $X_n$ arises from history $\gamma$, then
\[
  F(X_n) \subseteq \Delta(\Omega_0)
\]
is a face determined uniquely by $\gamma$.
\end{proposition}

Thus persistence corresponds to geometric stabilization rather than
substrate invariance.

\section{Minimal Commitment}
\label{sec:minimal}

Let $D$ be a target equivalence class. Define
\[
  \Gamma(D) = \{\gamma : \gamma(\Omega_0) \in D\}.
\]

\begin{definition}[Minimal commitment]
A history $\gamma \in \Gamma(D)$ is minimal if
\[
  S[\gamma] \le S[\gamma']
\]
for all $\gamma' \in \Gamma(D)$.
\end{definition}

\begin{proposition}
Minimal commitment histories maximize sheaf gluing prior to collapse.
\end{proposition}

\begin{proof}
Collapse reduces action by $\log|\ker q|$.
Gluing preserves entropy.
Hence minimal action is achieved by postponing collapse until necessary.
\end{proof}

\section{Finitude}
\label{sec:finitude}

\begin{theorem}[Finite descent]
Any strictly entropy-decreasing sequence in $\SP$ terminates.
\end{theorem}

\begin{proof}
Entropy is bounded below by zero and decreases strictly under Pop/Ref.
Hence no infinite strictly decreasing chain exists.
\end{proof}

\begin{corollary}
The deterministic fragment of $\SP$ is strongly normalizing.
\end{corollary}

\section{Variational Structure of Irreversibility}
\label{sec:variational}

Let $\gamma = (e_1,\dots,e_n)$ be a history in $\SP$.
Recall the action functional:
\[
  S[\gamma] = \sum_{i=1}^n \Delta\Opt_i
\]
where $\Delta\Opt_i = \Ent(\Omega_{i-1}) - \Ent(\Omega_i)$.

\begin{proposition}
$S[\gamma] = \Ent(\Omega_0) - \Ent(\Omega_n)$.
\end{proposition}

\begin{proof}
Telescoping sum.
\end{proof}

Thus the action depends only on endpoints.

\begin{definition}[Extremal history]
A history $\gamma$ between $X$ and $Y$ is extremal if
\[
  S[\gamma] = \min_{\gamma' : X \to Y} S[\gamma'].
\]
\end{definition}

\begin{theorem}[Endpoint minimality]
A history is extremal iff it performs no redundant collapse before
maximal gluing under admissibility constraints.
\end{theorem}

\begin{proof}
Collapse reduces entropy but may introduce unnecessary identification.
Since action depends only on entropy difference, any collapse that does
not change $\Ent$ must correspond to non-essential equivalence.
Such collapses can be postponed without altering endpoints.
\end{proof}

Hence irreversibility admits a variational characterization: the system
selects histories minimizing entropy expenditure subject to structural
constraints.

Under the realization functor $F$, this becomes:
\[
  S[\gamma] = \int_{M} \eta_\gamma \, d\mu.
\]
Thus discrete action equals integrated entropy slack in $\RSVP$.

\section{Adjoint Structure of Discrete--Continuous Realization}
\label{sec:adjunction}

Let $F : \SP \to \RSVP$ be the realization functor.
Define a coarse-graining functor
\[
  G : \RSVP \to \SP
\]
by mapping $(M,\Phi,v,S)$ to the partition of $M$ into
coherence basins under gradient flow of $\Phi$.

\begin{definition}
$G(M,\Phi,v,S)$ is the option-space whose elements are
stable attractor components of the flow induced by $v$.
\end{definition}

\begin{proposition}
$G$ is entropy non-increasing.
\end{proposition}

\begin{theorem}[Adjoint approximation]
There exists a natural transformation
\[
  \eta : \id_{\SP} \Rightarrow G \circ F
\]
which is initial among entropy-preserving embeddings.
\end{theorem}

\begin{proof}[Sketch]
Each $\Omega$ embeds into $\Delta(\Omega)$.
Coherence basins correspond to vertices.
The unit map identifies discrete elements with Dirac measures.
Universality follows from the simplex universal property.
\end{proof}

Thus discrete option-spaces are initial among coherence fields whose
basins reproduce their combinatorics.

While full adjunction may require restriction to a subcategory of
$\RSVP$, the construction demonstrates that the discrete--continuous
duality is structurally tight.

\section{Classification of Entropy-Monotone Irreversible Categories}
\label{sec:classification}

Let $\mathcal{C}$ be a symmetric monoidal irreversible category
equipped with a monotone resource functional $R$.

\begin{definition}
$\mathcal{C}$ is freely generated by a set $\mathcal{G}$ of primitives
if every morphism is a finite tensor-composition of elements of $\mathcal{G}$
subject only to resource monotonicity constraints.
\end{definition}

\begin{theorem}[Initiality of $\SP$]
Let $\mathcal{C}$ be a resource-monotone irreversible symmetric monoidal
category generated by four primitives satisfying:
\begin{enumerate}
  \item Strict entropy descent,
  \item Admissibility refinement without entropy change,
  \item Causal-preserving quotient,
  \item Annotated elimination.
\end{enumerate}
Then there exists a unique symmetric monoidal functor
\[
  H : \SP \to \mathcal{C}
\]
preserving generators and resource functional.
\end{theorem}

\begin{proof}
By free generation of $\SP$ under entropy axioms and universality of free symmetric monoidal categories.
\end{proof}

Hence $\SP$ is the initial object in the category of entropy-monotone
irreversible symmetric monoidal categories generated by four primitives.



\section{The Adjunction \texorpdfstring{$F\dashv G$}{F ⊣ G} via the Simplex-Realization Subcategory}
\label{sec:adjunction-precise}

Section~\ref{sec:adjunction} introduced a coarse-graining functor $G:\RSVP\to\SP$
defined by coherence basins of gradient flow, and asserted an adjoint approximation.
The obstruction to a clean adjunction on all of $\RSVP$ is that basin-selection is
not functorial under general smooth maps. This section repairs the statement by
restricting to a reflective subcategory $\RSVP_{\mathrm{simp}}\subseteq\RSVP$ on
which $G$ is defined precisely, and proves $G\circ F=\id_{\SP}$ on the nose.

\subsection{The Simplex-Realization Subcategory}

\begin{definition}[Realization fields]
\label{def:real-fields}
For each option-space $(\Omega,\mathcal{A})$, let $\mathcal{K}(\Omega,\mathcal{A})$
denote the set of pairs $(\Phi,v)$ in $C^\infty(\Delta(\Omega))\times\mathfrak{X}
(\Delta(\Omega))$ of the form produced by the object-construction of $F$
(Definition~\ref{def:F-objects}): $\Phi=\Phi_\Omega$ and $v=v_\Omega+\sum_{(i,j)\in B}
\delta v_{ij}$ for some finite binding set $B\subseteq\Omega^2$. Call such pairs
\emph{admissible realization fields}.
\end{definition}

\begin{definition}[Simplex-realization subcategory]
\label{def:rsvp-simp}
Define $\RSVP_{\mathrm{simp}}$ to be the full subcategory of $\RSVP$ spanned by
objects of the form $(\Delta(\Omega),\Phi,v,S_\Omega)$ where $(\Phi,v)\in
\mathcal{K}(\Omega,\mathcal{A})$ for some option-space $(\Omega,\mathcal{A})$,
with $S_\Omega$ the Shannon entropy density. Morphisms in $\RSVP_{\mathrm{simp}}$
are those in $\RSVP$ generated under composition and tensor by:
\begin{enumerate}[label=(\alph*)]
  \item face inclusions $(\iota_U,\eta_U)$ realizing $\Pop_U$ events;
  \item coarse-graining maps $(\varphi_\sim,\eta_\sim)$ realizing $\Col_\sim$ events;
  \item identity-on-manifold maps $(\id,0)$ with a single allowed $v$-increment
    $\delta v_{ij}$ realizing $\Bind_{ij}$ events.
\end{enumerate}
\end{definition}

By construction, $F$ lands in $\RSVP_{\mathrm{simp}}$: the image of every generating
morphism of $\SP$ is one of the generators (a), (b), (c) above.

\subsection{The Discretization Functor \texorpdfstring{$G$}{G}}

\begin{definition}[Discretization functor]
\label{def:G-functor}
Define $G:\RSVP_{\mathrm{simp}}\to\SP$ as follows.
\begin{itemize}
  \item \emph{On objects.} Send $(\Delta(\Omega),\Phi,v,S_\Omega)$ to $(\Omega,\mathcal{A})$,
    where $\mathcal{A}$ is read off from the binding structure recorded in $v$:
    specifically, $i\preceq j$ in $\mathcal{A}$ iff the field $v$ contains a
    $\delta v_{ij}$ increment.
  \item \emph{On generating morphisms.} Set $G(\iota_U,\eta_U):=\Pop_U$,
    $G(\varphi_\sim,\eta_\sim):=\Col_\sim$, and $G(\id,0\text{ with }\delta v_{ij})
    :=\Bind_{ij}$.
  \item \emph{On composites.} Extend by strict monoidal functoriality.
\end{itemize}
\end{definition}

\begin{proposition}[$G$ is a well-defined strict symmetric monoidal functor]
\label{prop:G-functor}
$G:\RSVP_{\mathrm{simp}}\to\SP$ is a well-defined strict symmetric monoidal functor.
\end{proposition}

\begin{proof}
On objects, $G$ is well-defined because each object of $\RSVP_{\mathrm{simp}}$ carries
a unique $(\Omega,\mathcal{A})$ by definition. On morphisms, $G$ sends generators to
generators by construction, and the defining relations of $\SP$
(Definition~\ref{def:free-gen}) are satisfied in $\RSVP_{\mathrm{simp}}$: pop-idempotence
corresponds to face-inclusion composition; bind-commutativity follows from the commutativity
of independent $v$-increments. Strict monoidality: $G(\Delta(\Omega_1)\times\Delta(\Omega_2))=
G(\Delta(\Omega_1\times\Omega_2))=(\Omega_1\times\Omega_2,\cdot)=(\Omega_1,\cdot)\otimes
(\Omega_2,\cdot)=G(\Delta(\Omega_1))\otimes G(\Delta(\Omega_2))$.
\end{proof}

\subsection{The Adjunction}

\begin{theorem}[$G\circ F=\id_{\SP}$ and the unit-counit pair]
\label{thm:adjunction}
\begin{enumerate}[label=(\roman*)]
  \item $G\circ F=\id_{\SP}$ strictly (not merely up to natural isomorphism).
  \item There is a natural transformation $\varepsilon:F\circ G\Rightarrow
    \id_{\RSVP_{\mathrm{simp}}}$ whose components are the identity morphisms
    on $\RSVP_{\mathrm{simp}}$-objects, making $F$ left adjoint to $G$.
  \item The adjunction unit $\eta:\id_{\SP}\Rightarrow G\circ F$ is the identity
    natural transformation.
\end{enumerate}
\end{theorem}

\begin{proof}
\textbf{(i).} On objects: $G(F(\Omega,\mathcal{A}))=G(\Delta(\Omega),\Phi_\Omega,v_\Omega,
S_\Omega)=(\Omega,\mathcal{A})$ by Definition~\ref{def:G-functor}. On generators:
$G(F(\Pop_U))=G(\iota_U,\eta_U)=\Pop_U$, and similarly for $\Bind$ and $\Col$.
Since both are strict monoidal functors and $\SP$ is freely generated, this forces
$G\circ F=\id_{\SP}$.

\textbf{(ii).} For each object $M=(\Delta(\Omega),\Phi,v,S_\Omega)\in\RSVP_{\mathrm{simp}}$,
define $\varepsilon_M:=\id_M$. This is a morphism in $\RSVP_{\mathrm{simp}}$ (the
identity). Naturality: for any morphism $h:M\to M'$ in $\RSVP_{\mathrm{simp}}$,
the square $\varepsilon_{M'}\circ FG(h)=h\circ\varepsilon_M$ reduces to
$\id_{M'}\circ h = h\circ\id_M$, which holds trivially.

\textbf{(iii).} The adjunction unit is $\eta_{(\Omega,\mathcal{A})}:=\id_{(\Omega,\mathcal{A})}$,
which is consistent with $G\circ F=\id_{\SP}$.
\end{proof}

\begin{remark}
The adjunction is tautological by design: $\RSVP_{\mathrm{simp}}$ is defined as the
image-presentation of $F$, so $G$ inverts $F$ on the nose. The non-trivial content
is (a) that the restriction $\RSVP_{\mathrm{simp}}\subseteq\RSVP$ is closed under
the operations used in Section~\ref{sec:functor}, which follows from
Definition~\ref{def:rsvp-simp}, and (b) that the asymmetry theorem
(Theorem~\ref{thm:asymmetry}) is preserved: the full $\RSVP$ strictly contains
$\RSVP_{\mathrm{simp}}$ (there exist smooth maps and smooth entropy densities not
of simplex form), so no extension of $G$ to all of $\RSVP$ making
$G\circ F=\id_{\SP}$ can be both total and faithful.
\end{remark}


% ═════════════════════════════════════════════════════════════════════════════
\section{Mathematical Status and Failure Modes}
\label{sec:status}

\subsection{Fully Formal}

The following are rigorous within standard mathematics:
(1) The category $\SP$ and its generating relations (Section~\ref{sec:sp});
(2) The causal preorder and the universal property of $\Col$ (Theorem~\ref{thm:collapse-universal});
(3) The strict symmetric monoidal structure (Proposition~\ref{prop:monoidal});
(4) Initiality of $\SP$ in $\EDSMC$ (Theorem~\ref{thm:initial});
(5) The meld theorem (Theorem~\ref{thm:sheafification});
(6) The $\RSVP$ category with witnessed morphisms (Proposition~\ref{prop:rsvp-cat});
(7) Functoriality of $F$ (Theorem~\ref{thm:functor});
(8) Strong normalization of the deterministic fragment (Theorem~\ref{thm:sn}).

\subsection{Modeling-Dependent}

The following require additional choices:
(1) The probability simplex $\Delta(\Omega)$ as the geometric realization of option-spaces
(alternatives: Wasserstein spaces, Hilbert spaces over $\Omega$);
(2) The specific entropy-transport PDE governing $\RSVP$ dynamics;
(3) The explicit form of $\delta\vF_{ij}$ in $F(\Bind)$.

\subsection{Failure Modes}

\begin{enumerate}[label=(\roman*)]
\item \textit{Collapse violating causality.} As demonstrated in
  Section~\ref{sec:worked}, applying $\Col_{a\sim b}$ after $\Bind_{ab}$ (which
  establishes $a\preceq b$ with $\downset{a}\ne\downset{b}$) violates causal
  admissibility. The system rejects such collapses in $\SP_c$.
\item \textit{Bind creating cycles.} A sequence $\Bind_{ab}\circ\Bind_{ba}$ creates
  $a\preceq b\preceq a$ with $a\ne b$, violating Lemma~\ref{lem:antisymmetry}. The
  well-formedness condition (acyclicity of the dependency graph) must be enforced
  externally or by a confluence condition.
\item \textit{Entropy functional failure.} If $p$ is not a probability measure on
  $\Omega$ (e.g., unnormalized), the Shannon entropy need not satisfy monotonicity
  under subsets. The framework requires that the entropy functional be consistently
  defined on all option-spaces with a fixed normalization.
\item \textit{$F$ failing to be faithful.} If two distinct binding constraints
  $\Bind_{ij}$ and $\Bind_{ij'}$ produce the same vector field modification
  $\delta\vF_{ij}=\delta\vF_{ij'}$ (because the coupling depends only on the simplex
  geometry, not on the labels), then $F$ is not faithful on $\Bind$ events. This
  requires the coupling to be injective in the label pair $(i,j)$, which is a
  non-trivial constraint on the realization.
\end{enumerate}

\subsection{Open Problems}

\begin{enumerate}
\item \textbf{Faithfulness of $F$.} Full characterization of when $F$ is faithful,
  particularly for $\Bind$ events.
\item \textbf{Entropy functional for $\Bind$.} Construction of a functional making
  Axiom~\ref{ax:entropy} strictly non-trivial for $\Bind$ (e.g., symmetry entropy
  $\mathsf{H}_{\mathrm{sym}}=\log|\Omega|-\log|\mathrm{Aut}(\Omega,\mathcal{A})|$).
\item \textbf{Compact-closed completion.} Whether $\SP$ admits a traced symmetric
  monoidal completion (the obstruction to full compact-closure is that $\eta:\unit\to
  X^*\tensor X$ would create optionality, violating Axiom~\ref{ax:irrev}).
\item \textbf{Almost-sure normalization} in the probabilistic fragment.
\item \textbf{$\infty$-categorical enhancement} replacing $\SP$ with an
  $(\infty,1)$-category where 2-morphisms are equivalences between histories.
\item \textbf{Derived Enhancement.}
Can $\SP$ be enhanced to an $(\infty,1)$-category where collapse corresponds to homotopy coequalizers rather than strict quotients?
\item \textbf{Entropy--Symmetry Duality.}
Construct an entropy functional on admissibility lattices via groupoid cardinality:
\[
  \mathsf{H}_{\mathrm{sym}}(\Omega) = \log |\mathrm{Aut}(\Omega)|.
\]
\item \textbf{Measure-Theoretic Completion.}
Replace $\Delta(\Omega)$ with the Wasserstein space $\mathcal{P}_2(\Omega)$ and study whether $F$ lifts to a functor into gradient-flow categories.
\item \textbf{Faithfulness Criterion.}
Characterize minimal geometric data on $\Delta(\Omega)$ making $F$ faithful. Conjecture: a non-degenerate Riemannian metric with curvature-sensitive slack is sufficient.
\item \textbf{Compact Completion and Trace.}
Determine whether adding formal time-reversal objects yields a traced monoidal category whose fixed points correspond to cyclic commitment structures.
\end{enumerate}

% ═════════════════════════════════════════════════════════════════════════════
\section{Dual Descriptions and Scale Separation}
\label{sec:dual}

The relationship between $\SP$ and $\RSVP$ is a scaling limit, not a symmetric duality.

\begin{center}
\renewcommand{\arraystretch}{1.4}
\begin{tabular}{ll}
\toprule
\textbf{$\SP$} & \textbf{$\RSVP$} \\
\midrule
Rewriting system & Thermodynamic limit \\
Discrete, combinatorial & Smooth, differential \\
Log-bound, append-only & Field-bound, diffusive \\
Constraint accumulation & Coherence redistribution \\
$\Col$ = coequalizer & $F(\Col)$ = renormalization \\
Action monotone increasing & PDE, parabolic \\
\bottomrule
\end{tabular}
\end{center}

Informally: as event step-size $\epsilon\to 0$ and event frequency $\to\infty$,
the discrete action $\mathcal{S}[\gamma]$ converges to a continuous action functional
whose Euler--Lagrange equation is the entropy-transport PDE. The functor $F$ is the
finite-step approximation. A rigorous derivation requires specifying convergence in a
Wasserstein topology on $\Delta(\Omega)$; this is recorded in
Appendix~\ref{app:scaling}.

% ═════════════════════════════════════════════════════════════════════════════
\section{Further Directions}
\label{sec:further}

\begin{enumerate}
\item \textbf{$\infty$-categorical enhancement.} Replace $\SP$ with an $(\infty,1)$-
  category in which 2-morphisms are policy-equivalences between histories and higher
  morphisms encode coherence conditions. Meld would then correspond to a homotopy pushout.
\item \textbf{Stochastic extension.} Incorporate a Markov chain on $\SP$ (transition
  probabilities on generators) and connect to the Independent Channels Lemma.
\item \textbf{Renormalization duality.} Formalize the duality between fine-grained
  $\SP$ categories and coarse-grained $\RSVP$ dynamics as a Wilsonian effective theory.
\item \textbf{Thermodynamic bounds.} Derive Landauer-type bounds on the minimum action
  required to realize a given collapse, using the discrete mechanics of
  Section~\ref{sec:mechanics}.
\item \textbf{Numerical simulation.} Discretize the entropy-transport PDE on $\Delta(\Omega)$
  and simulate convergence to coherence attractors for specific histories.
\end{enumerate}

% ═════════════════════════════════════════════════════════════════════════════
\section{Philosophical Remark}
\label{sec:philosophy}

Irreversibility in Spherepop is not an empirical claim about thermodynamic systems. It
is Axiom~\ref{ax:irrev}: the absence of inverses is constitutive of the framework. The
connection to Landauer's principle and Bennett's logical reversibility is structural
analogy, not reduction. This insulates the formalism from physical objections while
preserving the thermodynamic intuition.

The phrase ``joy of Spherepop'' names a phenomenological consequence of this structure:
in a universe of infinite options, nothing is decided. It is only by popping futures,
binding promises, and collapsing distinctions that a world becomes livable. The
functor $F$ makes this visible geometrically: each commitment sharpens the coherence
potential, directs the flow, and moves the probability mass toward a vertex.

% ═════════════════════════════════════════════════════════════════════════════
\section{Conclusion}
\label{sec:conclusion}

The paper compresses to three sentences.

\medskip
\noindent
\emph{Spherepop is the initial object in the category of entropy-decreasing symmetric
monoidal rewriting categories, freely generated by Pop, Refuse, Bind, and Collapse,
with worldhood as the sheaf condition on local event proposals and meld as
policy-induced sheafification.}

\medskip
\noindent
\emph{RSVP is a smooth entropy-witnessed field category whose morphisms carry
explicit slack data ensuring categorical closure without boundary-term ambiguity.}

\medskip
\noindent
\emph{$F:\SP\to\RSVP$ is a symmetric monoidal functor translating option-elimination
into boundary sharpening, dependency into directed coupling, and quotient collapse
into renormalization, with the structural asymmetry that $\SP$ accumulates constraint
while $\RSVP$ redistributes coherence---two scale-dual realizations of irreversible
entropy flow.}

\medskip
Constraint makes worlds. Entropy measures their cost. History defines identity.
The functor renders persistence geometrically visible.

% ═════════════════════════════════════════════════════════════════════════════
\appendix

\section{Confluence of Generator Relations}
\label{app:confluence}

The relations of Definition~\ref{def:free-gen} are confluent in the sense that any
two reduction paths starting from the same term reach the same normal form. This
follows from the fact that:
(a) idempotence of $\Pop$ is a local rewriting rule with no critical pairs;
(b) commutativity of independent $\Bind$ events ($\{i,j\}\cap\{k,l\}=\emptyset$)
is Church-Rosser by inspection (both orderings produce the same admissibility family
since the binding constraints are on disjoint pairs of elements);
(c) the monoidal laws are standard and their confluence is part of the theory of
symmetric monoidal categories.

\section{Free Generation: Formal Construction}
\label{app:free}

Let $\mathcal{G}$ be the directed multigraph:
\begin{itemize}
\item Vertices: all option-spaces $(\Omega,\mathcal{A})$.
\item Edges: for each nonempty $U\subsetneq\Omega$, edges $\Pop_U$ and $\RefOp_U$;
  for each pair $i\ne j\in\Omega$ with $\{i\preceq j\}\notin\mathcal{A}$, an edge
  $\Bind_{ij}$; for each causally admissible $\sim$, an edge $\Col_\sim$.
\end{itemize}
$\SP$ is the quotient of the free category $\mathcal{F}(\mathcal{G})$ by the
congruence of Definition~\ref{def:free-gen}, restricted to entropy-monotone paths.

\section{Symmetry Entropy for Bind}
\label{app:symm-entropy}

Define:
\[
  \mathsf{H}_{\mathrm{sym}}(\Omega,\mathcal{A})
  := \log|\Omega| - \log|\mathrm{Aut}(\Omega,\mathcal{A})|,
\]
where $\mathrm{Aut}(\Omega,\mathcal{A})$ is the group of admissibility-preserving
automorphisms of $(\Omega,\mathcal{A})$. Under $\Bind_{ij}$:
\begin{itemize}
\item $|\Omega|$ is unchanged, so $\log|\Omega|$ is unchanged;
\item $|\mathrm{Aut}(\Omega,\mathcal{A}')|<|\mathrm{Aut}(\Omega,\mathcal{A})|$ since
  the automorphism that swaps $i\leftrightarrow j$ (and is in $\mathrm{Aut}(\Omega,
  \mathcal{A})$ when $i,j$ are symmetric) is no longer admissibility-preserving after
  binding $i\preceq j$.
\end{itemize}
Hence $\mathsf{H}_{\mathrm{sym}}(\Omega,\mathcal{A}')>\mathsf{H}_{\mathrm{sym}}(\Omega,
\mathcal{A})$: symmetry entropy strictly increases under Bind.

This provides the entropy functional making Axiom~\ref{ax:entropy} non-trivially
active for $\Bind$. The monotonicity direction is reversed from the other generators
(symmetry entropy increases), which is consistent with $\Bind$ being entropy-preserving
in the Shannon sense while creating new asymmetry.

\section{Almost-Sure Normalization}
\label{app:asnorm}

Let Choice events assign probability $p_i$ to branch $i$. Define expected optionality
$\mathbb{E}[\Opt_{t+1}]=\sum_i p_i\Opt(X_i^{(t)})$. If there exists $\delta>0$ such
that $\mathbb{E}[\Opt_{t+1}]\le\Opt(X_t)-\delta$ whenever $\Opt(X_t)>0$, then
$(\Opt(X_t))_t$ is a non-negative supermartingale, hence $\Opt(X_t)\to 0$ almost
surely by Doob's convergence theorem. The condition $\delta>0$ holds when every
branch contains at least one $\Pop$ event with nonzero probability---a weak condition
satisfied by all well-typed Spherepop programs.

\section{Scaling Limit Heuristic}
\label{app:scaling}

Let $\epsilon>0$. Discretize time as $t_k=k\epsilon$. Suppose events occur at rate
$1/\epsilon$ and each $\Pop$ eliminates an $\epsilon$-fraction of mass. The discrete
action $\mathcal{S}[\gamma_\epsilon]=\sum_k L_k^{(\epsilon)}\approx\epsilon\sum_k
(-\dot\Opt(\gamma(t_k)))$. As $\epsilon\to 0$:
\[
  \mathcal{S}[\gamma_\epsilon] \to \int_0^T(-\dot\Opt(\gamma(t)))\,dt
  = \Opt(\Omega_0)-\Opt(\Omega_T).
\]
The Euler--Lagrange equation for this action functional under the constraint of
entropy-transport is $\partial_t\SF=\nabla\cdot(\kappa\nabla\PhiF)$, the RSVP PDE.
A rigorous derivation requires specifying the topology on the space of histories and
the sense of convergence; the natural setting is $\Gamma$-convergence of functionals
on $L^2(\Delta(\Omega))$.

\section{Faithfulness: Sufficient Conditions}
\label{app:faithful}

\begin{proposition}
$F$ is faithful on the subcategory generated by $\Pop$ and $\Col$ alone, provided the
Riemannian metric on $\Delta(\Omega)$ is non-degenerate and the boundary-sharpening
profile $\eta_U$ is injective as a function of $U$ (i.e., $U_1\ne U_2\Rightarrow
\eta_{U_1}\ne\eta_{U_2}$ as smooth functions on $\Delta(\Omega\setminus U_1)\cap
\Delta(\Omega\setminus U_2)$).
\end{proposition}

\begin{proof}
Two distinct $\Pop$ events $\Pop_{U_1}\ne\Pop_{U_2}$ produce distinct face inclusions
$\iota_{U_1}\ne\iota_{U_2}$ (they have different domains). Two distinct $\Col$ events
have different equivalence relations $\sim_1\ne\sim_2$, hence different coarse-graining
maps. The injectivity condition on $\eta$ ensures that even when the smooth maps $\varphi$
coincide (which happens only for trivial cases), the slack data distinguish the morphisms.
\end{proof}

\section{Notation Summary}
\label{app:notation}

\begin{center}
\renewcommand{\arraystretch}{1.4}
\begin{tabular}{ll}
\toprule
\textbf{Symbol} & \textbf{Meaning} \\
\midrule
$\SP$ & Free symmetric monoidal entropy-decreasing rewriting category \\
$\RSVP$ & Smooth entropy-witnessed field category \\
$\EDSMC$ & Category of entropy-decreasing symmetric monoidal categories \\
$(\Omega,\mathcal{A})$ & Option-space with admissibility family \\
$\Ent$ & Entropy functional $\Ob(\SP)\to\Rnn$ \\
$\Opt$ & Optionality functional $\Ob(\SP)\to\Rnn$ \\
$\Pop,\RefOp,\Bind,\Col$ & Four generating morphisms \\
$\Meld_\pi$ & Policy-induced sheafification operator \\
$\tensor$ & Symmetric monoidal tensor (Cartesian product) \\
$\unit$ & Unit object (trivial option-space) \\
$\preceq$ & Causal preorder on $\Omega$ \\
$\downset{x}$ & Causal past of $x$ \\
$\sim_q$ & Causally admissible equivalence (collapse policy) \\
$F:\SP\to\RSVP$ & Geometric realization functor \\
$(\varphi,\eta)$ & RSVP morphism with entropy-slack witness \\
$\Delta(\Omega)$ & Probability simplex over $\Omega$ \\
$\PhiF,\vF,\SF$ & Coherence potential, velocity field, entropy density \\
$\kappa$ & Diffusion coefficient in entropy-transport PDE \\
$\mathcal{S}[\gamma]$ & Action of history $\gamma$ \\
$\pi_t$ & Commitment (conjugate to optionality) \\
$H_t$ & Hamiltonian (remaining freedom) \\
$\mathcal{T}$ & Presheaf of local event proposals \\
$a_\pi(\mathcal{T})$ & Policy sheafification of $\mathcal{T}$ \\
$\eta_\pi$ & Universal $\pi$-invariant map \\
$\Kc$ & Kolmogorov complexity \\
$\mathcal{B}$ & Accounting functor tracking $\RefOp$ tags \\
\bottomrule
\end{tabular}
\end{center}

% ═════════════════════════════════════════════════════════════════════════════
\begin{thebibliography}{99}

\bibitem[Awodey(2010)]{awodey2010}
S.~Awodey, \emph{Category Theory}, 2nd ed., Oxford University Press, 2010.

\bibitem[Bennett(1973)]{bennett1973}
C.~H.~Bennett, Logical reversibility of computation,
\emph{IBM Journal of Research and Development}, 17(6):525--532, 1973.

\bibitem[Coquand \& Huet(1988)]{coquand1988}
T.~Coquand and G.~Huet, The calculus of constructions,
\emph{Proceedings of LICS}, 266--278, IEEE, 1988.

\bibitem[Fritz(2020)]{fritz2020}
T.~Fritz, A synthetic approach to Markov kernels, conditional independence,
and theorems on sufficient statistics,
\emph{Advances in Mathematics}, 370:107239, 2020.

\bibitem[Girard(1987)]{girard1987}
J.-Y.~Girard, Linear logic,
\emph{Theoretical Computer Science}, 50(1):1--101, 1987.

\bibitem[Giry(1982)]{giry1982}
M.~Giry, A categorical approach to probability theory,
\emph{Categorical Aspects of Topology and Analysis}, 68--85, Springer, 1982.

\bibitem[Green \& Altenkirch(2008)]{GreenAltenkirch2006}
A.~Green and T.~Altenkirch, From reversible to irreversible computations,
\emph{Electronic Notes in Theoretical Computer Science}, 210:65--74, 2008.
\newblock doi:\href{https://doi.org/10.1016/j.entcs.2008.04.018}{10.1016/j.entcs.2008.04.018}.

\bibitem[Grothendieck(1971)]{grothendieck1971}
A.~Grothendieck, \emph{S\'eminaire de G\'eom\'etrie Alg\'ebrique du Bois Marie (SGA 1)},
Springer, 1971.

\bibitem[Jacobson(1995)]{jacobson1995}
T.~Jacobson, Thermodynamics of spacetime: the Einstein equation of state,
\emph{Physical Review Letters}, 75(7):1260--1263, 1995.

\bibitem[Johnstone(2002)]{johnstone2002}
P.~T.~Johnstone, \emph{Sketches of an Elephant: A Topos Theory Compendium},
Oxford University Press, 2002.

\bibitem[Kolmogorov(1965)]{kolmogorov1965}
A.~N.~Kolmogorov, Three approaches to the quantitative definition of information,
\emph{Problems of Information Transmission}, 1(1):1--7, 1965.

\bibitem[Kozen(1981)]{kozen1981}
D.~Kozen, Semantics of probabilistic programs,
\emph{Journal of Computer and System Sciences}, 22(3):328--350, 1981.

\bibitem[Lambek \& Scott(1986)]{lambek1986}
J.~Lambek and P.~J.~Scott, \emph{Introduction to Higher Order Categorical Logic},
Cambridge University Press, 1986.

\bibitem[Landauer(1961)]{landauer1961}
R.~Landauer, Irreversibility and heat generation in the computing process,
\emph{IBM Journal of Research and Development}, 5(3):183--191, 1961.

\bibitem[Lawvere(1970)]{lawvere1970}
F.~W.~Lawvere, Quantifiers and sheaves,
\emph{Actes du Congr\`es International des Math\'ematiciens}, 1:329--334, 1970.

\bibitem[Mac Lane(1998)]{maclane1998}
S.~Mac Lane, \emph{Categories for the Working Mathematician}, 2nd ed., Springer, 1998.

\bibitem[Mac Lane \& Moerdijk(1992)]{maclane1992}
S.~Mac Lane and I.~Moerdijk, \emph{Sheaves in Geometry and Logic}, Springer, 1992.

\bibitem[Martin-L\"of(1975)]{martinlof1975}
P.~Martin-L\"of, An intuitionistic theory of types: predicative part,
\emph{Logic Colloquium '73}, 73--118, North-Holland, 1975.

\bibitem[Milner(1999)]{milner1999}
R.~Milner, \emph{Communicating and Mobile Systems: The $\pi$-Calculus},
Cambridge University Press, 1999.

\bibitem[Moggi(1991)]{moggi1991}
E.~Moggi, Notions of computation and monads,
\emph{Information and Computation}, 93(1):55--92, 1991.

\bibitem[Pantev et al.(2013)]{pantev2013}
T.~Pantev, B.~To\"en, M.~Vaqui\'e, and G.~Vezzosi, Shifted symplectic structures,
\emph{Publications Math\'ematiques de l'IH\'ES}, 117(1):271--328, 2013.

\bibitem[Pierce(2002)]{pierce2002}
B.~C.~Pierce, \emph{Types and Programming Languages}, MIT Press, 2002.

\bibitem[Plotkin(2004)]{plotkin2004}
G.~D.~Plotkin, A structural approach to operational semantics,
\emph{Journal of Logic and Algebraic Programming}, 60--61:17--139, 2004.

\bibitem[Shannon(1948)]{shannon1948}
C.~E.~Shannon, A mathematical theory of communication,
\emph{Bell System Technical Journal}, 27:379--423, 1948.

\bibitem[Street(1972)]{street1972}
R.~Street, The formal theory of monads,
\emph{Journal of Pure and Applied Algebra}, 2(2):149--168, 1972.

\bibitem[Verlinde(2011)]{verlinde2011}
E.~Verlinde, On the origin of gravity and the laws of Newton,
\emph{Journal of High Energy Physics}, 2011(4):29, 2011.

\bibitem[Wilson(1983)]{wilson1983}
K.~G.~Wilson, The renormalization group and critical phenomena,
\emph{Reviews of Modern Physics}, 55(3):583--600, 1983.

\end{thebibliography}

\end{document}